% REVISION CHANGELOG (in response to GENETICS-2026-309068 reviewer reports)
%
% ABSTRACT
%   - Rewritten to reflect the new BV-to-fitness scope (matches the
%     introduction's keystone framing).
%   - Now states the R^2 identity for Ag and the tight 1 - Ag bound
%     explicitly.
%   - Names the regime explanation (high-Ag elite breeding, low-Ag
%     near-optimum natural populations) and the counterintuitive
%     prediction.
%   - Mentions the Lande-Arnold dictionary connection.
%   - Removed tentative phrases ("explores one possible explanation",
%     "may frequently reside") and the generic closing.
%
% ARTICLE SUMMARY
%   - Lightly tightened to add the breeding-vs-natural contrast.
%
% INTRODUCTION
%   - Replaced with a four-paragraph scope statement.
%   - Para 1 distinguishes the G-to-P map question from the BV-to-fitness
%     question; takes breeding values as primitive (Reviewer 1, Editor).
%   - Para 2 commits explicitly to the BV-to-fitness question.
%   - Para 3 lists four deliverables, balanced for breeders AND
%     evolutionary biologists; names Lande/Schluter/Fisher and the
%     Hansen-Houle-style locality conditions explicitly.
%   - Para 4 adds the bridge observation: even on a perfectly additive
%     G-to-P map, a quadratic fitness function generates fitness epistasis;
%     mechanistic vs. geometric non-additivity are distinct.
%   - Removed the additivity-vs-interactions schematic figure (now done
%     in prose).
%
% NEW SUBSECTION: Section 2.4 "Notation: relation to Lande-Arnold (1983)"
%   - Direct response to Reviewer 2 Concern #2.
%   - Explicit prose mapping of (b, H, G) to (beta, gamma, G, P, E).
%   - Resolves the W(a) vs W(z) confusion via Equation eq:Wa_definition
%     and the marginalisation-over-e identity.
%   - Phillips-Arnold relation H = gamma - beta*beta^T derived in S7.
%   - Environmental buffering H_a vs H_z derived in S7.
%   - Table 1 (notation) was expanded to absorb the dictionary content,
%     so we deliver the dictionary without adding a second table
%     (Reviewer 1 had asked for fewer tables; we end at 2, down from
%     the original 4).
%
% REWORKED: Section 3.6 "Generalised eigenvalue rotation and the
%   basis-free form of Ag"
%   - Direct response to the editor's request for an explicit
%     generalised eigenvalue rotation.
%   - Whitening (Delta a = G^{1/2} u) and rotation (orthogonal U
%     diagonalising H-tilde = G^{1/2} H G^{1/2}) make the rotation
%     transparent; mu_i are eigenvalues of HG = generalised eigenvalues
%     of (H, G^{-1}).
%   - Basis-free form Ag = sum c_i^2 / (sum c_i^2 + 0.5 sum mu_i^2).
%   - Scale-freeness now stated as a property with derivation in
%     Appendix S1 (new subsection "Coordinate invariance of V_lin
%     and V_quad").
%   - Schluter line of least resistance interpretation enriched with
%     the complementary "small-|mu|" axis (direction along which
%     standing variance feels least curvature).
%
% NEW APPENDIX: S7 "Notation Dictionary Derivations"
%   - H = gamma - beta*beta^T, with derivation at the population mean.
%   - H_a = -(Omega+E)^{-1} via Gaussian marginalisation, linking the
%     dictionary to the Level 2 selection update.
%   - Step-by-step recipe for computing Ag from a published LA study.
%
% APPENDIX S1 ADDITION
%   - New subsection "Coordinate invariance of V_lin and V_quad" with
%     the chain-rule and trace-cyclicity proof, supporting the
%     scale-freeness claim in Section 3.6.
%
% NEW PARAGRAPH: Section 4 "Beyond the scalar closure: a directional
%   retention operator"
%   - Acknowledges the scalar rho closure as a known simplification.
%   - Names the directional retention operator rho_i as the formal
%     object that would replace it, varying along the principal
%     directions of HG (using the rotation language from Section 3.6).
%   - Identifies three primitives that would generate rho_i:
%     recombination rates, locus-by-trait architecture, per-direction
%     selection strength.
%   - Bounds the deferral by appealing to the framework's deliberate
%     agnosticism about locus-level architecture, while noting that
%     the diagnostic Ag itself does not depend on rho.
%   - Inserted between the equilibrium-variance discussion (Eq 22)
%     and the equilibrium-Ag subsection so the caveat covers both
%     analyses.
%
% DISCUSSION
%   - Replaced "Connections to natural populations" with one
%     mutation/selection paragraph and one substantive deployment
%     paragraph (four uses for evolutionary biologists, parallel to
%     the breeding subsection).
%   - Removed the asexual E. coli (Khan 2011) worked example, both the
%     Discussion subsection and Figure: TODO marker preserves position
%     for the Nature-Plants-derived diploid replacement (decided with MC).
%   - Tightened "What the framework does not offer" so it doesn't
%     duplicate the introduction's scope statement.
%   - Tightened "Empirical challenges and a measurement pathway" by
%     merging the two paragraphs that both discussed DMS.
%
% DIMENSIONAL CONSIDERATIONS
%   - Consolidated 8 paragraphs into 4. Removed redundant restatement
%     of the effective-dimension argument (Reviewer 1, p. 13-14
%     duplication).
%   - Removed fig:keff (figure now lives in SI; section reads cleanly
%     in prose).
%
% FIGURES AND TABLES (main text)
%   - Figures: 11 -> 5. Removed: additivity_vs_interactions, fig:phase
%     (heuristic regime map), fig:keff, fig:khan_worked_example,
%     fig:neighbourhoods (point carried by fig:channel and fig:framework
%     B/C panels), fig:Ag (sigma^2 scaling fully explicit in eq:Vlin,
%     eq:Vquad and the crossover algebra).
%   - The five surviving figures map cleanly to the paper's distinct
%     argumentative jobs: fig:channel (concept), fig:Ag_validation
%     (empirical evidence), fig:framework (mechanism + regime snapshot),
%     fig:adaptation-trajectory (rise-then-fall dynamic), fig:trajectories
%     (long-run regime contrast).
%   - Tables:  4 -> 2. Removed: tab:Ag_interpretation (overlaps prose),
%     tab:regimes (overlaps fig:trajectories), tab:natural_vs_breeding
%     (in FAQ; redundant). Table 1 (tab:notation) was expanded to
%     include LA correspondences so a separate dictionary table
%     was unnecessary -- direct response to Reviewer 1 (fewer tables)
%     AND Reviewer 2 (notation dictionary) in one move.
%
% APPENDIX S6 (FAQ) -- now S8 after the renumbering, see below
%   - Q13 cleaned: removed placeholder cite tags, removed duplicate
%     table, removed the redundant non-specialist summary at the end
%     (article summary at the top already serves this role).
%
% MISC
%   - Added missing \label{sec:level1} that an Appendix S3 ref pointed
%     to (pre-existing bug; no longer unresolved).
%   - All begin/end pairs balance; all internal refs resolve.
%
% APPLIED: Tier 3/4 surgical fixes (subset)
%   - C^2 regularity statement: added "Provided ell is at least twice
%     continuously differentiable in a neighbourhood of a_0" to the
%     Taylor-setup sentence in Section 2.2.
%   - a_0 = E[a] clarification: changed "typically the population
%     mean" to the explicit definition "a_0 = E[a]" with the centering
%     consequence "E[Delta a] = 0 by construction" stated. Updated
%     the downstream "key point deserves emphasis" paragraph for
%     consistency.
%
% APPLIED: Cooper-Technow co-origin and dimension-enlargement scope
%   - Discussion / "Connections to breeding programs": added a short
%     paragraph crediting Cooper for first observing that the rise of
%     %GCA under K=15 isolation in the E(NK) simulations of Technow et
%     al. (2021) resembled the variance dynamics of the framework, and
%     stating the geometric reading: conversion of variance under drift
%     equals contraction of G along principal directions of HG, with
%     per-direction rate set by variance retention. Inserted between
%     paragraphs 1 and 2 of the breeding subsection.
%   - Discussion / "Connections to natural populations": added a
%     scoping paragraph after the fourth use restricting the rotating-b
%     interpretation to environmental shifts that re-weight existing
%     selective pressures. Names dimension-enlarging shifts (novel
%     pathogen, unfamiliar climate, new herbicide) as outside the
%     local-geometric framework as currently formulated. Triplet
%     intentionally pairs natural-population and breeding-program
%     examples in line with the paper's communion-of-fields posture.
%
% APPLIED: New Appendix S5 -- Stochastic-simulation validation
%   (originally drafted as S8 and then renumbered to S5 to group with
%   Appendix S4, the other empirical illustration; see Appendix
%   renumbering note below.)
%   - Added a full new appendix (now S5) reporting SLiM forward-time
%     individual-based simulation validation of the framework. The
%     appendix has six subsections covering simulation design, the
%     A_g = R^2 identity across stochastic realisations (Fig S5.1),
%     across-replicate variability of framework invariants vs their
%     constituent eigenvalues (Fig S5.2), geometric invariants k_eff
%     and condition(G) of the genetic covariance (Fig S5.3), and a
%     synthesis with explicit limitations and BV-level scope statement.
%   - Three figures are inserted (multiseed_4trait_moving.pdf,
%     cv_4trait_moving.pdf, geom_invariants_4trait_moving.pdf) using
%     updated Appendix S5 simulation outputs and a local figure-counter
%     renumbering so they print as S5.1--S5.3 rather than continuing
%     the main-text count.
%   - Three new bibkeys are introduced and used only in S5:
%     HallerMesser2019, Walsh2018, Stern2013. The other cites
%     (CarterHermissonHansen2005, Hansen2013, Schluter1996,
%     CheverudRoutman1995, Hansen2006) reuse existing keys.
%   - Pre-empts the Carter et al. (2005) concern raised by R1 by
%     explicitly distinguishing variance-redistribution (what we
%     observe) from variance-expansion (what Carter et al. predict
%     under nonlinear G->P maps), and by scoping the validation
%     explicitly to the BV-level dynamics that the framework treats
%     (the G->P map is upstream of where the framework operates).
%
% APPLIED: Appendix renumbering for conceptual coherence
%   The appendices have been reordered so that math derivations come
%   first (S1 Variance Identities, S2 Selection Operators, S3
%   Equilibrium Analysis), the worked example follows (S4), the two
%   empirical appendices are adjacent (new S5 Simulation Validation
%   sits next to S4 Worked Example), the math notation bridge
%   follows (new S6 Notation Dictionary Derivations, formerly S7),
%   and discursive material is at the end (new S7 Historical Context,
%   formerly S5 Intellectual Lineage; new S8 FAQ, formerly S6).
%   This affected: four section headings, one inline cross-reference
%   in the main text (Appendix S6 FAQ -> Appendix S8 FAQ), the
%   figure counter and label naming inside the simulation appendix
%   (S8.x -> S5.x; fig:s8_* -> fig:s5_*), and the comment block
%   inside the simulation appendix that referred to it as the "last
%   appendix."  Internal LaTeX \ref{} cross-references resolve via
%   labels rather than numbers and required no changes.
%
% APPLIED: Rho derivations in Appendix S2 (editor request)
%   - Added new subsection at the end of Appendix S2 titled "The LD
%     retention coefficient rho: derivations under the Lande, Bulmer,
%     and Barton-Turelli frameworks." Derives rho = 1/2 in the Lande
%     infinitesimal limit under random mating, rho = E[1-r_ij] in the
%     Bulmer model with non-uniform recombination, rho = (1+F)/2 under
%     partial selfing (translating the verbal "inbreeding raises rho"
%     into a formal monotone-in-F statement), and a leading eigenvalue
%     of the recombination operator in the Barton-Turelli framework.
%     Closes with a summary table in prose.
%   - Three new bibkeys introduced: Wright1933, HaymanMather1953,
%     Bulmer1980. The other cites (BartonTurelli1991, Burger2000,
%     Bulmer1971, Wright1969) reuse existing keys.
%   - Directly addresses the editor's specific request: "show what rho
%     is then: e.g. in Lande model or Bulmer model or Weissman-Barton
%     model... how inbreeding in breeding programs makes rho higher."
%
% APPLIED: Anisotropy paragraph after Eq 28 (editor request)
%   - Added explicit paragraph "On the asymmetric anisotropy in this
%     argument" right after the participation-ratio Eq 28. Explains
%     that the isotropic-H treatment in Eqs 27-28 is a deliberate
%     didactic choice to isolate the dimensional structure of G from
%     that of H, points the reader to Section 3.6 (the rotation
%     section, sec:rotation) for the fully anisotropic case via
%     simultaneous diagonalisation, and connects to Technow et al.
%     (2026) as an empirical instance of an anisotropic and
%     history-dependent H emerging from E(NK) simulations.
%   - One new bibkey: Technow2026.
%
% APPLIED: Appendix S7 title change
%   - Renamed "Appendix S7: Intellectual Lineage" to "Appendix S7:
%     Historical Context" per author preference. Section content
%     unchanged; label app:lineage retained for backwards
%     compatibility with internal cross-references.
%
% APPLIED: Four post-audit textual additions (May 2026)
%   - Section 3.1 (Variance from the linear term): Hansen-Houle
%     evolvability sentence appended after the local-form Fisher's
%     fundamental theorem statement, identifying V_lin = b' G b as
%     e(b), so the connection to standard evolvability terminology
%     appears at the introduction of V_lin rather than only in the
%     Discussion. Bibkey: HansenHoule2008.
%   - Section 3.2 (Variance from the quadratic term): opening
%     sentence expanded with the multivariate-normality contrast.
%     Var(Q) depends on fourth moments of Delta a and therefore
%     requires the Gaussian assumption, unlike Var(L) = V_lin which
%     only requires a centred distribution with covariance G. The
%     asymmetry between Eq Vlin and Eq Vquad is now flagged at the
%     equations rather than only in Appendix S1.
%   - Section 4 (Level 1 derivation), sentence "Selection compresses
%     variance by adding precision.": citation cluster appended.
%     Bibkeys: Lande1980, Bulmer1971, Bulmer1980, BartonTurelli1991.
%     Anchors the precision-addition reading at first invocation.
%   - Section 7 (Discussion), tangent-plane sentence: expanded with
%     citations to Rice (2008) and Milocco & Salazar-Ciudad (2022)
%     for the nonlinear-G-to-P-map case, plus a closing sentence
%     stating that the contribution of A_g is to make the locality
%     condition quantitative. Bibkeys: Rice2008,
%     MiloccoSalazarCiudad2022.
% =====================================================================

\documentclass[9pt,twocolumn,twoside,lineno]{gsajnl}

\usepackage{epstopdf}
\articletype{inv}

\runningtitle{Additive Channels in Fitness Landscapes}
\runningauthor{Ortiz-Barrientos \& Cooper}

\usepackage{amsmath,amssymb,amsthm,bm}
\usepackage{natbib}
% \usepackage{nameref}  % not used: section cross-references are written explicitly because gsajnl can leave \nameref blank
\usepackage{graphicx}
\usepackage{float}
\usepackage{caption}
\graphicspath{{./additiveChannels/figures/}{./}}
\usepackage{booktabs,array}
\usepackage{microtype}
\usepackage{xspace}
\emergencystretch=2em
\allowdisplaybreaks

% ----------------------------------------------------------
% CUSTOM COMMANDS
% ----------------------------------------------------------
\newcommand{\vect}[1]{\boldsymbol{#1}}
\newcommand{\mat}[1]{\mathbf{#1}}
\newcommand{\RR}{\ensuremath{\mathbb{R}}\xspace}
\newcommand{\E}{\ensuremath{\mathbb{E}}\xspace}
\newcommand{\tr}{\ensuremath{\operatorname{tr}}\xspace}
\newcommand{\Ag}{\ensuremath{\mathcal{A}_{g}}\xspace}
\newcommand{\Gsel}{\ensuremath{\mat{G}^{\text{sel}}}\xspace}
\newcommand{\Ggenic}{\ensuremath{\mat{G}_{\text{genic}}}\xspace}
\newcommand{\Vlin}{\ensuremath{V_{\text{lin}}}\xspace}
\newcommand{\Vquad}{\ensuremath{V_{\text{quad}}}\xspace}
\newcommand{\keff}{\ensuremath{k_{\mathrm{eff}}}\xspace}
\newcommand{\Var}{\ensuremath{\operatorname{Var}}\xspace}
\newcommand{\Cov}{\ensuremath{\operatorname{Cov}}\xspace}
\newcommand{\Corr}{\ensuremath{\operatorname{Corr}}\xspace}


% ----------------------------------------------------------
% TITLE & AUTHORS
% ----------------------------------------------------------
\title{Additive Channels in Curved Fitness Landscapes}

\author[1,2,*]{Daniel Ortiz-Barrientos}
\author[2,3]{Mark Cooper}

\affil[1]{School of the Environment, The University of Queensland, St Lucia, QLD 4072, Australia}
\affil[2]{ARC Centre for Plant Success in Nature and Agriculture, The University of Queensland, St Lucia, QLD 4072, Australia}
\affil[3]{Queensland Alliance for Agriculture and Food Innovation, The University of Queensland, St Lucia, QLD 4072, Australia}

\affil[*]{Correspondence should be addressed to d.ortizbarrientos@uq.edu.au and mark.cooper@uq.edu.au}

\date{\today}

% ----------------------------------------------------------
% ABSTRACT & KEYWORDS
% ----------------------------------------------------------
\begin{abstract}
Additive models of inheritance predict the short-term response to selection remarkably well, even when the underlying biology involves widespread dominance and gene--gene interaction. We argue that this success reflects a property of how fitness varies with the additive genetic component of a trait, not a property of the molecular architecture beneath. We make this quantitative through an additivity index, \Ag, that measures the fraction of local log-fitness variance captured by a linear approximation, with the remainder attributable to local curvature of the fitness surface. Under Gaussian/quadratic assumptions, \Ag equals the squared correlation between the linear approximation and local log-fitness, so $1-\Ag$ is the fraction of local log-fitness variance not captured by a purely linear predictor. We call regions of breeding-value space where \Ag is high \emph{additive channels} and develop a coupled selection--inheritance framework that identifies when populations enter, persist in, and leave them. The framework predicts that stable, well-adapted populations can be those for which additive prediction is least informative for directional-response prediction.
\end{abstract}

\keywords{fitness landscapes; quantitative genetics; epistasis; additive variance; response to selection}

% ----------------------------------------------------------
\begin{document}
% ----------------------------------------------------------

\maketitle
\thispagestyle{firststyle}
\vspace{-13pt}

% ----------------------------------------------------------
% ARTICLE SUMMARY (GENETICS requirement: ~100 words for non-specialists)
% ----------------------------------------------------------
\noindent\textbf{Article summary:} Gene interactions are common, yet additive genetic models often predict short-term evolution. We propose that additivity is a local property of where populations sit in breeding-value space on a curved fitness landscape, not a global property of the biology beneath. We define an index, \Ag, that compares slope variance to curvature variance and identifies populations in additive channels for which linear prediction performs well. The framework predicts that elite breeding pools enter such channels under sustained selection, while natural populations near a fitness optimum may leave them because the directional signal collapses.

\vspace{6pt}

% ==========================================================
\section{Introduction}
\label{sec:intro}
% ==========================================================

Quantitative genetics asks several distinct questions about additivity, and the literature does not always distinguish them. One question concerns the genotype-to-phenotype map. Do allelic effects on a trait combine additively, or do dominance and epistasis dominate? Decades of work have shown that average effects remain well-defined and often predictively useful even when the underlying mechanisms are highly nonlinear \citep{Fisher1918, CheverudRoutman1995, Hansen2006, Hansen2013, Rice2002, CarterHermissonHansen2005, Morrissey2015, GonzalezForero2026}. This is a question about how genotypes encode traits.

A second question is geometric. Given a population whose standing variation lives in a small neighbourhood of breeding-value space, when does the tangent plane of log-fitness $\ell(\vect{a})$ capture most of the variance that selection acts on? Here we are no longer asking how genotypes become traits. We are asking how a population, statistically described by a mean breeding value $\bar{\vect{a}}$ and a covariance matrix $\mat{G}$, samples the curvature of a fitness surface. This is the question we address. We take breeding values as primitive; Fisher's least-squares decomposition guarantees their existence regardless of how nonlinear the genotype-to-phenotype map may be, and the geometric question that follows can be answered without committing to a particular mechanistic model of gene interaction.

Why should this question matter? First, it gives a single quantity, the additivity index \Ag, computable from $\vect{b}$, $\mat{H}$, and $\mat{G}$, that diagnoses whether linear prediction should work in a given population, turning the locality conditions under which Lande's equation, Schluter's lines of least resistance, and Fisher's fundamental theorem operate into something one can compute. Second, it places a tight upper bound on the gain from any non-additive extension, whether epistatic terms in genomic prediction or quadratic $\bm{\gamma}$ terms in selection-gradient regression. Third, it explains why populations occupy the regimes they do, elite breeding pools settle into high-\Ag regions where sustained directional selection holds $\vect{b}$ away from zero, whereas natural populations near a fitness optimum drift toward low-\Ag regions where $\vect{b} \to \vect{0}$ collapses $\Vlin$ without collapsing $\Vquad$, implying counterintuitively that the stable, well-adapted populations most often studied in selection-gradient work may be those where additive prediction is least reliable. And fourth, it reconciles two empirical patterns that have long appeared in tension: pervasive molecular epistasis on the one hand \citep{Phillips2008, HuangMackay2016} and the predictive success of additive models on the other \citep{HillGoddardVisscher2008, Meuwissen2001}. Both can be true once we recognise that additivity is not a global property of the genetic architecture but a \emph{local} property of where the population sits on a curved surface. See Appendix~\ref{app:lineage} for historical antecedents to our work.

% ----------------------------------------------------------
% TABLE 1: NOTATION + LANDE-ARNOLD DICTIONARY (merged)
% -----------------------------------------------------------
\begin{table}[H]
\centering
\caption{Notation used throughout the manuscript and its relation to \citet{LandeArnold1983} (LA). LA work in phenotype space and report regression coefficients of \emph{relative} fitness; we work in breeding-value space with the local Taylor coefficients of \emph{log} fitness. The \textit{Notation: relation to Lande--Arnold (1983)} subsection explains the three differences (domain, scale, environmental buffering) and Appendix~\ref{app:dictionary} derives the relations.}
\label{tab:notation}
\small
\begin{tabular}{@{}l p{1.9cm} p{4.0cm}@{}}
\toprule
Symbol & Meaning & Notes \\
\midrule
$\vect{a}$ & Breeding values & Additive genetic values; we work in BV space, LA in phenotype space \\
$\vect{z}$ & Phenotype & $\vect{z} = \vect{a} + \vect{e}$; identical to LA's $\vect{z}$ \\
$\vect{e}$ & Env.\ deviation & $\vect{e} \sim \mathcal{N}(\vect{0}, \mat{E})$ \\
$\mat{G}$ & Genetic covariance & $\mat{G} = \Var(\vect{a})$; identical to LA's $\mat{G}$ \\
$\mat{E}$ & Env.\ covariance & $\mat{E} = \Var(\vect{e})$; identical to LA's $\mat{E}$ \\
$\mat{P}$ & Phenotypic cov. & $\mat{P} = \mat{G} + \mat{E}$; identical to LA's $\mat{P}$ \\
$W(\vect{a})$ & Fitness in BV space & $\mathbb{E}_{\vect{e}}[W(\vect{a}+\vect{e})]$ (Eq.~\ref{eq:Wa_definition}); LA's $W(\vect{z})$ is pointwise \\
$\ell = \log W(\vect{a})$ & Log-fitness & Local geometry object; LA work with relative $w = W/\bar{W}$ \\
$\vect{b}$ & Log-fitness gradient & $\nabla_{\vect{a}} \ell$; drives $\Vlin$; $\vect{b}\approx\bm{\beta}$ to leading order under weak selection and after mapping to the same domain (see footnote~\ref{fn:b-vs-beta}) \\
$\mat{H}$ & Log-fitness Hessian & $\nabla^2_{\vect{a}} \ell$; drives $\Vquad$; $\mat{H} = \bm{\gamma}-\bm{\beta}\bm{\beta}^{\top}$ in same domain (Eq.~\ref{eq:H_vs_gamma}) \\
$\Vlin$ & Linear variance & $\vect{b}^\top \mat{G} \vect{b}$; equals Hansen--Houle evolvability $e(\vect{b})$ \\
$\Vquad$ & Quadratic variance & $\tfrac{1}{2}\tr[(\mat{H}\mat{G})^2]$; no standard LA name \\
$\Ag$ & Additivity index & $\Vlin/(\Vlin + \Vquad)$; introduced here \\
$L$ & Linear Taylor term & $L = \vect{b}^\top \Delta\vect{a}$; $\Var(L) = \Vlin$ (Eq.~\ref{eq:Vlin}) \\
$Q$ & Quadratic Taylor term & $Q = \tfrac{1}{2}\Delta\vect{a}^\top \mat{H} \Delta\vect{a}$; $\Var(Q) = \Vquad$ (Eq.~\ref{eq:Vquad}) \\
$R$ & Curvature relevance ratio & $R = \Vquad/\Vlin$; small means linear prediction wins (Eq.~\ref{eq:R-isotropic}) \\
$\Gamma$ & Curvature magnitude & $\Gamma = -\mat{H}$ (pos.\ def.); selection width $\Omega = \Gamma^{-1}$ \\
$\rho$ & LD retention & Mating-system parameter \\
$\mat{M}$ & Mutational input & Variance per generation \\
$\Ggenic$ & Genic variance & From allele frequencies \\
$\bm{\beta}$ & LA linear gradient & $\mat{P}^{-1}\vect{S}$; phenotype space \\
$\vect{S}$ & Selection differential & $\vect{S} = \mat{P}\bm{\beta}$; phenotypic mean shift after selection (LA) \\
$\bm{\gamma}$ & LA quadratic gradient & Phenotype space; relates to $\mat{H}$ via Eq.~\ref{eq:H_vs_gamma} \\
$\keff$ & Effective dimension & $(\tr \mat{G})^2 / \tr(\mat{G}^2)$ \\
$\sigma^2$ & Variance scale & Typical breeding-value variance; $\sigma^2 = \tr(\mat{G})/n$ in isotropic case \\
$\delta$ & Displacement from optimum & $\delta = |\bar{\vect{a}} - \vect{a}^*|$; controls $\vect{b} = -\Gamma\bm{\delta}$ near a quadratic optimum \\
\bottomrule
\end{tabular}
\end{table}

One observation makes the geometric question worth taking seriously on its own terms. Even if the genotype-to-phenotype map were perfectly additive, every allele contributing independently with no dominance and no epistasis, the resulting log-fitness surface in breeding-value space need not be. A quadratic fitness function applied to additive breeding values produces log-fitness terms that are, by construction, non-additive in those breeding values. Stabilising selection on additive traits generates fitness epistasis even in the absence of any mechanistic gene interaction. The two sources of fitness non-additivity, the \emph{mechanistic} (from gene interactions in the G--P map) and the \emph{geometric} (from curvature in the BV--fitness map), are distinct phenomena that can be studied separately. The framework developed here isolates the second. We call regions where the geometric source contributes little to heritable fitness variation \emph{additive channels}, and we ask when populations enter them, persist in them, and leave them.

We develop the framework through three components, illustrated in Figure~\ref{fig:channel}. The first concerns the local geometry of log-fitness surfaces: the slope (gradient $\vect{b}$) and curvature (Hessian $\mat{H}$) that a population experiences at its current location.\footnote{\label{fn:b-vs-beta}Two gradient vectors appear in this paper. The vector $\vect{b} = \nabla \ell |_{\vect{a}_0}$ is the gradient of log-fitness in BV-space at the reference point. The vector $\vect{\beta} = \mat{P}^{-1}\vect{S}$ is the Lande--Arnold selection gradient in phenotype-space. Under weak selection and Gaussian assumptions, after mapping both to the same domain, $\vect{\beta} \approx \vect{b}$ to leading order; outside this regime they are formally distinct objects.} The second concerns selection: how Gaussian selection operators transform the genetic covariance matrix $\mat{G}$ into a post-selection covariance $\Gsel$ within a single generation. The third concerns inheritance: how recombination, mutation, and linkage disequilibrium transmit $\Gsel$ across generations. The additivity index \Ag sits atop this framework, evaluating whether the resulting $\mat{G}$ places the population in a regime where linear prediction is adequate.
% ----------------------------------------------------------
% FIGURE 1: CHANNEL CONCEPT
% ----------------------------------------------------------
\begin{figure}[!htb]
\centering
\includegraphics[width=\columnwidth]{figure1_channel.pdf}
\caption{The additive channel concept. A population (blue ellipses) traverses a rugged fitness landscape along a relatively flat region. The vertical axis shows log-fitness $\ell(\vect{a}) = \log W(\vect{a})$, the natural scale for Gaussian selection models and the quantity whose local slope ($\vect{b}$) and curvature ($\mat{H}$) define the additivity index. As selection compresses genetic variance (shrinking ellipses), the population samples progressively smaller neighbourhoods where curvature contributes less to fitness variation. The trajectory (amber line) shows evolutionary movement through the channel. On a globally curved landscape, a population can behave additively if its variance is small enough that the quadratic term in the local Taylor expansion remains negligible.}
\label{fig:channel}
\end{figure}

We are explicit about scope. The framework treats breeding values as primitive and the fitness landscape as given, without modelling its emergence from molecular mechanisms. The Gaussian closures we use for tractability may break down under strong selection or skewed distributions \citep{TurelliBarton1994}, and the scalar parameter $\rho$ summarising LD retention is a coarse stand-in for a locus-specific process \citep{Walsh2018}. We view this paper as a diagnostic framework for when additive models should work, not as a complete theory of genetic architecture.

Different analyses in the paper hold different quantities constant, and we flag this here because the same variables play different roles in each. In the \textit{Local Geometry of Log-Fitness Surfaces} and \textit{The Additivity Index} sections, where we set up the local geometry and define $\Ag$, the gradient $\vect{b}$ and Hessian $\mat{H}$ at a fixed reference point are taken as given; the diagnostic acts pointwise on the triple $(\vect{b}, \mat{H}, \mat{G})$. In the \textit{Selection Reshapes Genetic Covariances} section, where Gaussian selection operators transform $\mat{G}$ within a generation, $\mat{H}$ (equivalently $\Gamma = -\mat{H}$) is held constant during the within-generation update; this is exact when the fitness surface is exactly log-quadratic, and approximate otherwise. In the \textit{Inheritance and Equilibrium} section, when we couple the selection operator to a Bulmer-style across-generation recursion, $\mat{H}$, the mutational input $\mat{M}$, and the LD-retention parameter $\rho$ are all held constant across generations; this corresponds biologically to a stationary fitness landscape, an unchanging mutational regime, and an unchanging mating system. The equilibrium analysis (the \textit{Inheritance and Equilibrium} section) holds these three constant in $\mat{G}$-space and asks where $\mat{G}$ settles. Throughout, the optimum is taken to be either fixed (natural-population analyses) or moving with sustained external pressure (breeding-program analyses); we make the choice explicit in each case. The framework therefore does not claim that $\mat{H}$ is unchanging in nature---it provides a local diagnostic that is most informative when the local landscape geometry is approximately stable on the timescale of the analysis.

Notation is summarised in Table~\ref{tab:notation}. We proceed by setting up the local geometry of log-fitness in the \textit{Local Geometry of Log-Fitness Surfaces} section, defining \Ag and its connection to linear predictability in the \textit{Additivity Index} section, deriving Gaussian selection operators in the \textit{Selection Reshapes Genetic Covariances} section, coupling them to a minimal across-generation inheritance recursion in the \textit{Inheritance and Equilibrium} section, and addressing dimensionality in the \textit{Dimensional Considerations} section. A worked numerical example that traces each operation back to a biological interpretation is given in Appendix~\ref{app:example}.

% ==========================================================
\section{Local Geometry of Log-Fitness Surfaces}
\label{sec:geometry}
% ==========================================================

The starting point of our analysis is the observation that any smooth function can be locally approximated by a polynomial. Applied to fitness, this means that the fitness surface in the neighbourhood of a population can be described by its slope (first derivatives) and curvature (second derivatives), regardless of how complicated the global surface might be. We work with log-fitness rather than fitness itself. This choice has three advantages: log-fitness is the natural scale for Gaussian selection, because a quadratic log-fitness surface implies fitness weights of the form $\exp(\,\vect{b}^\top\Delta\vect{a}-\tfrac{1}{2}\Delta\vect{a}^\top\Gamma\Delta\vect{a}\,)$ (Eq.~\eqref{eq:W-Gaussian}); multiplicative fitness effects become additive on the log scale, which aligns with common biological intuitions about independent contributions \citep{Rice2004}; and it avoids issues with fitness values near zero. Throughout, we write $\ell(\vect{a}) = \log W(\vect{a})$ for the log-fitness surface. We refer to this log-fitness representation together with the local Taylor approximation and its slope and curvature objects ($\vect{b}$ and $\mat{H}$) as the \emph{local-geometry setup}.

\subsection{The population as a cloud in breeding-value space}

Consider a population of individuals, each characterised by a vector of breeding values $\vect{a} \in \RR^n$. The breeding value is the additive genetic value, the portion of an individual's genotype that is transmitted to offspring and contributes additively to phenotype \citep{Fisher1918, Lynch1998}. For our purposes, the key point is that breeding values live in a continuous space, and a population is not a single point but a cloud of points in this space. This cloud has a center, which we can take to be the mean breeding value $\bar{\vect{a}}$. The cloud also has a shape, described by the additive genetic covariance matrix $\mat{G} = \Var(\vect{a})$. If $\mat{G}$ is large, the population spreads widely in breeding-value space; if $\mat{G}$ is small, the population is concentrated near its mean. The eigenvectors of $\mat{G}$ describe the principal directions of variation, and the eigenvalues define how much variance lies along each direction.

Now consider a log-fitness surface $\ell(\vect{a})$ defined over breeding-value space. This surface assigns a log-fitness value to each possible breeding value. The surface can be arbitrarily complicated in its global shape, with multiple peaks, ridges, valleys, and saddles \citep{Gavrilets2004}. But the population does not experience the entire surface. It experiences only the local neighbourhood that its genetic variation allows it to sample.

The typical parameter values describing this cloud differ substantially between natural and breeding contexts. In natural populations, environmental variance $\mat{E}$ is often large relative to genetic variance $\mat{G}$, so heritability $h^2 = G/(G + E)$ is moderate \citep{Visscher2008}. Mutation continuously adds genetic variance at a rate $\mat{M}$ per generation \citep{Lynch1998}. Recombination in outcrossing species erodes linkage disequilibrium rapidly, so the LD-retention parameter $\rho$ is typically around 0.5 or lower. 

In breeding programs, the picture differs \citep{Allard1999, FalconerMackay1996}. Environmental variance can be reduced through controlled trials and replication, raising heritability. Mutational input is negligible on the 10--20 generation timescales of typical breeding cycles. Inbreeding and selfing increase $\rho$, allowing selection-induced LD to persist \citep{Bulmer1971}. These differences have important consequences for where populations sit along the additive channel, because they control how strongly the Gaussian selection operator compresses $\mat{G}$ within generations and how the inheritance recursion restores (or fails to restore) variance across generations.

\subsection{The Taylor expansion: separating slope from curvature}

To make the local approximation precise, we expand log-fitness as a Taylor series around a reference point $\vect{a}_0$, taken throughout this paper to be the population mean, $\vect{a}_0 = \E[\vect{a}]$. Let $\Delta\vect{a} = \vect{a} - \vect{a}_0$ denote the deviation of an individual's breeding value from this reference, so that $\E[\Delta\vect{a}] = \vect{0}$ by construction. Provided $\ell$ is at least twice continuously differentiable in a neighbourhood of $\vect{a}_0$, the second-order Taylor expansion is:
\begin{equation}
\ell(\vect{a}) \approx \ell_0 + \vect{b}^\top \Delta\vect{a} + \tfrac{1}{2} \Delta\vect{a}^\top \mat{H} \Delta\vect{a}
\label{eq:taylor}
\end{equation}
Here $\ell_0 = \ell(\vect{a}_0)$ is log-fitness at the reference point, $\vect{b} = \nabla\ell|_{\vect{a}_0}$ is the gradient (a vector of first partial derivatives), and $\mat{H} = \nabla^2\ell|_{\vect{a}_0}$ is the Hessian matrix (a matrix of second partial derivatives).

Each term in this expansion has a clear interpretation. The constant $\ell_0$ sets the baseline log-fitness at the reference point. The linear term $\vect{b}^\top \Delta\vect{a}$ describes how log-fitness changes as we move away from the reference in any direction; it is the dot product of the gradient with the displacement. The quadratic term $\tfrac{1}{2}\Delta\vect{a}^\top \mat{H} \Delta\vect{a}$ captures the curvature of the surface; it quantifies how the linear approximation breaks down. The gradient $\vect{b}$ is the mathematical representation of directional selection \citep{LandeArnold1983}. 

If $\vect{b} = \vect{0}$, there is no local tendency for log-fitness to change in any direction, which occurs at fitness optima, minima, and saddle points. If $\vect{b} \neq \vect{0}$, selection favours movement in the direction of $\vect{b}$. The Hessian $\mat{H}$ is the mathematical representation of how selection changes as the population moves. If $\mat{H}$ is negative definite (all eigenvalues negative), the surface curves downward in all directions, indicating a fitness peak with stabilising selection. If $\mat{H}$ has mixed signs, the surface is saddle-shaped. The magnitude of the eigenvalues indicates how sharply the surface curves. This connects to the $\gamma$ matrix of nonlinear selection gradients \citep{LandeArnold1983, PhillipsArnold1989}.

A key point deserves emphasis: $\vect{b}$ and $\mat{H}$ are evaluated at the current population mean $\vect{a}_0$. If the mean moves, these quantities change even if the fitness surface itself is fixed. The additivity index we construct below is therefore a local property of ``population $\times$ landscape at that location,'' not a property of the landscape alone.

\subsection{When does the Taylor approximation hold?}

The Taylor expansion is valid when the population's genetic spread is small relative to the scale over which the fitness surface changes. More precisely, two conditions must hold. First, log-fitness must be approximately log-quadratic locally: the third- and higher-order terms of the Taylor expansion of $\ell(\vect{a})$ around $\vect{a}_0$ must be negligible across the range the population samples. Second, if $\sigma$ is a typical standard deviation of breeding values (say, $\sigma = \sqrt{\tr(\mat{G})/n}$), and $L$ is the length scale over which the Hessian changes appreciably, then the second-order approximation requires $\sigma \ll L$. The first condition is about the local shape of $\ell$; the second is about the spread of the population that samples it. Both are more likely to hold when genetic variance is small, precisely the regime that breeding programs tend to create through repeated selection. They may fail when genetic variance is large, when the fitness surface has sharp features, or when the population spans multiple fitness peaks. We revisit these limitations when we describe scope, empirical estimation, and failure modes of the Gaussian/quadratic closure.

\subsection{Notation: relation to Lande--Arnold (1983)}
\label{sec:notation_dictionary}

Two conventions for analysing local selection on quantitative traits coexist in the literature, and a reader fluent in one easily mistakes the symbols of the other. \citet{LandeArnold1983} work in phenotype space, regress \emph{relative} fitness on phenotype, and report a directional gradient $\bm{\beta}$ and a quadratic gradient $\bm{\gamma}$. The framework developed here works in breeding-value space, expands \emph{log} fitness around the population mean, and reports the gradient $\vect{b}$ and Hessian $\mat{H}$ of $\ell(\vect{a}) = \log W(\vect{a})$. The two conventions describe the same biology and reduce to one another under standard Gaussian assumptions, but they differ in three small ways that matter for any reader who wants to compute $\Vlin$, $\Vquad$, or \Ag from published Lande--Arnold estimates.

The first difference concerns the meaning of the fitness function. Lande--Arnold's $W(\vect{z})$ assigns a fitness to each phenotype; an individual's fitness is read off the surface at its phenotypic location. Our $W(\vect{a})$ assigns a fitness to each breeding value, defined as the expected fitness of an individual carrying that breeding value, marginalised over the environmental contribution $\vect{e}\sim\mathcal{N}(\vect{0},\mat{E})$:
\begin{equation}
W(\vect{a}) \;=\; \mathbb{E}_{\vect{e}}\!\left[W(\vect{a}+\vect{e})\right].
\label{eq:Wa_definition}
\end{equation}
The two are different functions even when both are Gaussian: marginalising over $\vect{e}$ smooths the surface seen in breeding-value space relative to the surface seen in phenotype space. For Gaussian stabilising selection of width $\bm{\Omega}$, the curvature satisfies $\mat{H}_{\vect{a}} = -(\bm{\Omega}+\mat{E})^{-1}$ in breeding-value space versus $\mat{H}_{\vect{z}} = -\bm{\Omega}^{-1}$ in phenotype space; environmental variance uniformly buffers the apparent curvature. This is the same buffering that appears in the Level~2 selection update on $\mat{G}$ (the \textit{Selection Reshapes Genetic Covariances} section); the dictionary makes the connection explicit.

The second difference concerns the scale on which fitness is regressed. Lande--Arnold use the relative fitness $w(\vect{z}) = W(\vect{z})/\bar{W}$ and report the regression coefficients $\bm{\beta}$ (linear) and $\bm{\gamma}$ (quadratic); we use $\ell(\vect{a}) = \log W(\vect{a})$ and report the local Taylor coefficients $\vect{b}$ and $\mat{H}$. Linear gradients agree to leading order, $\bm{\beta} \approx \vect{b}$. Curvature, however, transforms non-trivially. Evaluated in the same domain at the population mean where $w=1$, a short calculation (Appendix~\ref{app:dictionary}) gives
\begin{equation}
\mat{H} \;=\; \bm{\gamma}-\bm{\beta}\bm{\beta}^{\top}.
\label{eq:H_vs_gamma}
\end{equation}
The $\bm{\beta}\bm{\beta}^{\top}$ correction is the relation noted by \citet{PhillipsArnold1989}: principal-axis decomposition of $\bm{\gamma}$ alone does not recover the principal directions of stabilising selection on log fitness, because the directional component contributes to $\bm{\gamma}$ as well. Working with $\mat{H}$ rather than with $\bm{\gamma}$ removes this contamination by construction.

The third difference concerns the domain in which curvature is reported. In our framework $\mat{H}$ lives in breeding-value space, where $\Vquad = \tfrac{1}{2}\tr[(\mat{H}\mat{G})^2]$ has its cleanest interpretation as the variance contributed by a quadratic regression of log fitness on breeding values. Lande--Arnold's $\bm{\gamma}$ lives in phenotype space; recovering the BV-space $\mat{H}$ from a published $\bm{\gamma}$ requires both the Phillips--Arnold subtraction in Equation~\eqref{eq:H_vs_gamma} and the environmental-buffering correction. Table~\ref{tab:notation} collects the correspondences alongside the in-paper meaning of each symbol.

The practical implication is that any study reporting $(\bm{\beta},\bm{\gamma})$ together with estimates of $\mat{G}$, $\mat{E}$, and $\mat{P}$ contains, under the Gaussian assumptions used throughout this paper, the ingredients for $\vect{b}$, $\mat{H}$, $\Vlin$, $\Vquad$, and \Ag. The recipe is to compute $\mat{H} = \bm{\gamma}-\bm{\beta}\bm{\beta}^{\top}$ at the population mean, apply the environmental-buffering correction to recover the BV-space curvature, and substitute. We rely on this fact in the Discussion to argue that decades of selection-gradient estimates already contain the ingredients for the diagnostic, with no new data type required.

% ----------------------------------------------------------
% FIGURE fig:neighbourhoods: REMOVED in this revision.
% Its point (same G, different landscape locations -> different Ag) is
% carried by fig:channel (the population moves through the channel) and
% by fig:framework panels B/C (two populations, two Ag values).
% ----------------------------------------------------------

% ==========================================================
\section{The Additivity Index}
\label{sec:additivity}
% ==========================================================

We now construct a diagnostic that quantifies the extent to which the genetic variance in log-fitness arises from the linear versus the quadratic term in Equation~\eqref{eq:taylor}. This diagnostic, which we call the additivity index and denote by \Ag, provides a way to assess whether a population is in a regime in which additive prediction should work. We refer to the resulting partition into $\Vlin$, $\Vquad$, and their ratio \Ag as the \emph{variance decomposition}.

\subsection{Variance from the linear term}

Consider the linear term in the Taylor expansion: $L = \vect{b}^\top \Delta\vect{a}$. This is a random variable because $\Delta\vect{a}$ varies across individuals in the population. Its variance is:
\begin{equation}
\Vlin = \Var(\vect{b}^\top \Delta\vect{a}) = \vect{b}^\top \mat{G} \vect{b}
\label{eq:Vlin}
\end{equation}
This formula follows from a standard result: the variance of a linear combination of random variables is a quadratic form in the coefficients. A step-by-step derivation is given in Appendix~\ref{app:variance}. The quantity $\Vlin = \vect{b}^\top \mat{G} \vect{b}$ measures how much genetic variance lies along the gradient direction. If $\mat{G}$ has most of its variance along directions perpendicular to $\vect{b}$, then $\Vlin$ is small even if total genetic variance is large. This is the geometric content of Fisher's fundamental theorem in its local form \citep{Fisher1930, Price1972}: the rate of change of mean fitness depends not on total genetic variance but on genetic variance in the fitness direction. Equivalently, $\Vlin = \vect{b}^{\top}\mat{G}\vect{b}$ is the evolvability metric $e(\vect{b})$ of \citet{HansenHoule2008}, evaluated along the local log-fitness gradient: the magnitude of additive variance in the direction selection actually pushes.

\subsection{Variance from the quadratic term}

Consider the quadratic term: $Q = \tfrac{1}{2}\Delta\vect{a}^\top \mat{H} \Delta\vect{a}$. Its variance depends on fourth moments of $\Delta\vect{a}$ and therefore, unlike the linear-form variance $\Vlin$ in Equation~\eqref{eq:Vlin}, requires the additional assumption that $\Delta\vect{a}$ is multivariate normal. Computing it uses standard identities for quadratic forms of multivariate normal vectors \citep{MathaiProvost1992} (a derivation is given in Appendix~\ref{app:variance}):
\begin{equation}
\begin{split}
\Vquad = \Var(Q) &= \tfrac{1}{2}\tr(\mat{H}\mat{G}\mat{H}\mat{G})
= \tfrac{1}{2}\tr\!\left[\big(\mat{G}^{1/2}\mat{H}\mat{G}^{1/2}\big)^2\right] \\
&= \tfrac{1}{2}\big\|\mat{G}^{1/2}\mat{H}\mat{G}^{1/2}\big\|_F^2
\end{split}
\label{eq:Vquad}
\end{equation}
where $\|\cdot\|_F$ denotes the Frobenius norm. The final form makes non-negativity explicit and emphasises that $\Vquad$ depends on how the population covariance $\mat{G}$ weights curvature directions of $\mat{H}$.

This formula has a scaling structure that follows directly from the Taylor expansion itself. If we write $\mat{G} = \sigma^2 \mat{C}$ where $\mat{C}$ is a correlation matrix (trace $n$), then $\Vlin \propto \sigma^2$ while $\Vquad \propto \sigma^4$. This is a direct mathematical consequence of $L$ being linear and $Q$ quadratic in $\Delta\vect{a}$: the variance of a linear form scales as the second moment of its argument, while the variance of a quadratic form scales as the fourth. The biological implication is therefore equally direct: when genetic variance is reduced---by selection, drift, or programmatic compression---the curvature contribution shrinks faster than the linear contribution, so populations move into the additive channel quadratically faster than they lose evolvability.

In the isotropic case $\mat{G}=\sigma^2\mat{I}$ and $\mat{H}=-\gamma\mat{I}$,
the crossover variance where $\Vlin=\Vquad$ occurs at
\[
\sigma_\star^2=\frac{2|\vect{b}|^2}{n\gamma^2},
\]
which implies $\Ag(\sigma_\star^2)=0.5$. More generally, any threshold
$\Ag>\alpha$ corresponds to an upper bound on variance,
\[
\sigma^2 < \frac{(1-\alpha)}{\alpha}\cdot \frac{2|\vect{b}|^2}{n\gamma^2}.
\]
These expressions make clear that the boundary of the additive channel is set by the ratio of directional slope strength to curvature strength, scaled by the effective dimension $\keff$ (Table~\ref{tab:notation}; formal treatment in the \textit{Dimensional Considerations} section).

\subsection{The additivity index defined}

Before defining \Ag, we state a general identity. The total variance in log-fitness from the Taylor expansion is:
\begin{equation}
\Var(L + Q) = \Vlin + \Vquad + 2\Cov(L, Q)
\label{eq:total-var}
\end{equation}
Under multivariate normality, $\Cov(L, Q) = 0$ because all odd central moments of a Gaussian vanish (Appendix~\ref{app:variance}). This orthogonality justifies treating the linear and quadratic variances as additive when computing total log-fitness variance. If the distribution of breeding values deviates substantially from normality, this orthogonality may break down, and the decomposition becomes less clean.

Given Gaussian breeding values, we define the additivity index as the fraction of genetic log-fitness variance attributable to the linear term:
\begin{equation}
\Ag = \frac{\Vlin}{\Vlin + \Vquad} = \frac{1}{1 + R}
\label{eq:Ag}
\end{equation}
where $R = \Vquad/\Vlin$ is the curvature relevance ratio.

When $\Ag$ is close to 1, nearly all genetic log-fitness variance comes from the linear term. In this regime, the quadratic contribution is small, and linear prediction should work well. When $\Ag$ is close to 0, the quadratic term dominates, and linear prediction will miss most of the fitness variation among genotypes.

\subsection{\Ag as a local measure of linear predictability}
\label{sec:Ag_predictability}

So far, we have defined \Ag as a variance ratio. Next, we make explicit why this quantity is directly connected to the typical claim that ``linear (additive) prediction works.''

From the local quadratic approximation (Eq.~\eqref{eq:taylor}), deviations in log-fitness decompose as
$\tilde{\ell}=L+Q$, where $L=\vect{b}^\top\Delta\vect{a}$ is the linear term and $Q = \tfrac{1}{2}\Delta\vect{a}^\top \mat{H} \Delta\vect{a}$ is the quadratic term. Under the Gaussian assumption on $\Delta\vect{a}$, we have $\Cov(L,Q) = 0$.

The total variance of local log-fitness is therefore $\Var(\tilde{\ell}) = \Var(L) + \Var(Q) = \Vlin + \Vquad$. The covariance between the linear predictor and the true local log-fitness is $\Cov(L, \tilde{\ell}) = \Cov(L, L+Q) = \Var(L) + 0 = \Vlin$. The Pearson correlation between the linear predictor $L$ and the local log-fitness deviation $\tilde{\ell}$ is then:
\begin{equation}
\begin{split}
\Corr(L, \tilde{\ell}) &= \frac{\Cov(L, \tilde{\ell})}{\sqrt{\Var(L) \cdot \Var(\tilde{\ell})}} = \frac{\Vlin}{\sqrt{\Vlin \cdot (\Vlin + \Vquad)}} \\[4pt]
&= \sqrt{\frac{\Vlin}{\Vlin + \Vquad}} = \sqrt{\Ag}
\end{split}
\end{equation}
Equivalently,
\begin{equation}
\Corr(L, \tilde{\ell})^2 = \Ag
\label{eq:Ag_as_R2}
\end{equation}

This result shows that, within the local Gaussian/quadratic setting, \Ag is exactly the fraction of variance in local log-fitness that is explained by the linear term---an $R^2$-like measure of how well the best available linear predictor captures fitness differences across the breeding values the population samples. In genomic prediction, this corresponds to the expected accuracy of ranking individuals by estimated breeding values when the model ignores curvature, a situation typical of additive BLUP and genomic BLUP approaches.

\subsubsection*{Simulation validation}

Figure~\ref{fig:Ag_validation} tests two claims at once: that the identity $\Ag = R^2$ holds across realistic genetic architectures, and that the operational content of the additive channel---linear-only weighting reproduces full-quadratic selection outcomes when $\Ag$ is high---also holds across those architectures. We swept 360 random local-geometry triples $(\vect{b}, \mat{H}, \mat{G})$ spanning dimensions $n \in \{2, 4, 8\}$, both stabilising and saddle-shaped curvatures, and $\mat{G}$-magnitudes covering two orders of magnitude. For each triple, breeding values were sampled ($N = 20{,}000$) from three biologically realistic distributions: a Gaussian baseline; a polygenic finite-locus model (200 small-effect loci at intermediate frequency $p = 0.5$, near-Gaussian by the central limit theorem); and a rare-large-effect model (one locus at $p = 0.02$ with effect size ten times the polygenic baseline, plus the same 200-locus background). The first represents the classical idealisation; the second represents the dominant biological case (many loci of small effect, segregating at moderate frequency); the third represents introgression and the deliberate movement of new alleles into a breeding pool, where a single large-effect variant sits on top of an otherwise polygenic background.

Panel~(A) tests the $\Ag = R^2$ identity. All three architectures lie on the identity line $y = x$: median $|\Ag - R^2| < 0.007$ in every case. The polygenic and rare-large-effect cases track the Gaussian baseline because the central limit theorem applied to a finite mixture of bounded contributions produces near-Gaussian breeding-value distributions even when the underlying loci carry uneven weight. Heavy tails in the strict sense---segregating variance dominated by a small number of large-effect alleles at low frequency---can arise under introgression or migration, but our simulation shows they do not break the identity in the regime that matters for prediction.

Panel~(B) tests the operational claim. We compute the distance between the selected mean under linear-only weighting $\exp(L)$ and under full weighting $\exp(L+Q)$, normalised by a population-scale baseline. The pooled median (green line) drops by approximately twentyfold across the $\Ag$ range---from $\sim 40\%$ at $\Ag < 0.25$ to $\sim 2\%$ at $\Ag > 0.75$---and the drop is uniform across all three architectures. When the population sits inside the additive channel, additive approximations are difficult to distinguish from the full quadratic model, regardless of whether the breeding-value distribution is Gaussian, polygenic with intermediate-frequency loci, or contaminated by a rare large-effect variant. A more aggressive heavy-tailed stress test (multivariate $t_4$, with no genetic-mechanism analogue) is provided in the SI; the framework's predictions degrade there in the way the algebra predicts, but $t_4$ violates the central-limit assumptions that finite-locus polygenic and rare-large-effect architectures both satisfy, so we treat it as a mathematical worst case rather than a biological model.

% ----------------------------------------------------------
% FIGURE fig:Ag: REMOVED in this revision.
% Its three panels visualised the sigma^2 scaling (V_lin propto sigma^2,
% V_quad propto sigma^4, the crossover at sigma_star^2, and the resulting
% Ag decline). All of this is explicit in Equations eq:Vlin, eq:Vquad,
% and the surrounding crossover formulae, so the figure visualised but
% added no evidence beyond what the algebra already showed.
% ----------------------------------------------------------

% ----------------------------------------------------------
% FIGURE: A_g VALIDATION
% ----------------------------------------------------------
\begin{figure*}[!htb]
\centering
\includegraphics[width=\textwidth]{figure2_validation.pdf}
\caption{\textbf{The additivity index is robust to genetic architecture.} We tested the identity $\Ag = R^2$ and the operational claim that linear-only weighting reproduces full-quadratic selection outcomes across 360 random local-geometry triples $(\vect{b}, \mat{H}, \mat{G})$ spanning dimensions $n \in \{2, 4, 8\}$, both stabilising and saddle-shaped curvatures, and $\mat{G}$-magnitudes covering two orders of magnitude. For each triple, breeding values were sampled ($N = 20{,}000$) from three biologically realistic distributions: a Gaussian baseline; a polygenic finite-locus model (200 small-effect loci at intermediate frequency, $p = 0.5$, near-Gaussian by the central limit theorem); and a rare-large-effect model (one locus at $p = 0.02$ with effect size ten times the polygenic baseline, plus the same 200-locus background). \textbf{(A)}~Empirical squared correlation $R^2$ between the linear predictor $L = \vect{b}^\top \Delta\vect{a}$ (Eq.~\ref{eq:Vlin}) and the true local log-fitness $\tilde{\ell} = L + Q$ (Eqs.~\ref{eq:taylor},~\ref{eq:Vquad}), plotted against analytical $\Ag$ (Eq.~\ref{eq:Ag}). All three architectures lie on the identity line $y = x$: median $|\Ag - R^2| < 0.007$ in every case. The shaded region marks the additive channel ($\Ag \geq 0.75$). \textbf{(B)}~Relative difference between selected means under linear-only weighting $\exp(L)$ and full weighting $\exp(L + Q)$, normalised by $(\|\Delta\bar{\vect{a}}_{\mathrm{full}}\| + \sqrt{\tr(\mat{G})})$, on a log scale. The pooled median (green line) drops by approximately twentyfold across the $\Ag$ range---from $\sim 40\%$ at $\Ag < 0.25$ to $\sim 2\%$ at $\Ag > 0.75$---uniformly across all three architectures, showing that the operational claim does not depend on architecture either. The framework's predictions are therefore robust to the kinds of departures from Gaussianity that real genetic architectures produce. A more aggressive heavy-tailed stress test (multivariate $t_4$, with no genetic-mechanism analogue) is provided in the SI.}
\label{fig:Ag_validation}
\end{figure*}

\subsection{Interpretation and additional caveats}

Several additional caveats bear mention beyond those discussed above. At a fitness optimum where $\vect{b} = \vect{0}$, the linear variance $\Vlin = 0$. The quadratic variance $\Vquad = \tfrac{1}{2}\tr[(\mat{H}\mat{G})^2]$ remains strictly positive whenever the Hessian is non-zero, so $\Ag = 0/(\Vlin + \Vquad) = 0$ cleanly---there is no $0/0$ ambiguity at the optimum, only a vanishing numerator over a strictly positive denominator. This does not mean the genetics is ``non-additive'' or that ``epistasis dominates,'' it simply means there is no slope term to compare against curvature. The additive channel concept is therefore most informative when populations are displaced from an optimum, whether tracking a moving target, experiencing directional selection, or otherwise having $\vect{b} \neq \vect{0}$. This is an important diagnostic for off-optimum dynamics. A related limit case deserves a brief note. \Ag is built from a second-order Taylor expansion; at an inflection point of $\ell$, the local Hessian $\mat{H}$ vanishes in some directions, which sends $\Vquad$ along those directions to zero even when $\ell$ is not locally linear. This limit case matters because higher-order terms, rather than quadratic curvature, become the dominant source of nonlinearity in such regions, and \Ag will report the population as residing inside an additive channel even though linear prediction will degrade once the population traverses through the inflection. \Ag is therefore a diagnostic of \emph{quadratic-closure} additivity, not of additivity in any deeper sense; populations near inflection points are a known regime where the diagnostic can be locally optimistic. We discuss this regime further in Appendix~\ref{app:faq}. 

A second consideration concerns the Gaussian assumption discussed above. If the distribution of breeding values deviates substantially from normality, the orthogonality between $L$ and $Q$ may break down, and \Ag may not cleanly separate linear from quadratic contributions. Finally, because our application of the Taylor expansion truncates at second order, \Ag may mischaracterise systems in which the fitness surface exhibits substantial higher-order structure and genetic variance is large enough to probe it. These caveats reinforce the view that \Ag is best understood as a diagnostic tool rather than a complete description of genetic architecture.

\subsection{Generalised eigenvalue rotation and the basis-free form of \Ag}
\label{sec:rotation}

Three quantities control the additivity index: the gradient $\vect{b}$, the curvature $\mat{H}$, and the genetic covariance $\mat{G}$. The simplest case is isotropic, with $\mat{G} = \sigma^2 \mat{I}$ and $\mat{H} = -\gamma \mat{I}$, giving
\begin{align}
\Vlin &= |\vect{b}|^2 \sigma^2, \\
\Vquad &= \tfrac{1}{2} n \gamma^2 \sigma^4,
\end{align}
and the curvature relevance ratio
\begin{equation}
R = \frac{\Vquad}{\Vlin} = \frac{n \gamma^2 \sigma^2}{2 |\vect{b}|^2}.
\label{eq:R-isotropic}
\end{equation}
$R$ is small---and \Ag close to 1---when directional selection is strong, curvature is weak, variance is small, or the effective dimension is low. In the isotropic case the direction of $\vect{b}$ is irrelevant: every direction of breeding-value space is geometrically equivalent.

Real biology is anisotropic. Variance is concentrated along the principal axes of $\mat{G}$; curvature is concentrated along the principal axes of $\mat{H}$; the two need not coincide; and $\vect{b}$ has its own orientation. Two transformations expose the joint structure. \emph{Whitening} the breeding-value distribution via the substitution $\Delta\vect{a} = \mat{G}^{1/2}\vect{u}$, where $\mat{G}^{1/2}$ is the symmetric positive-definite square root, standardises the covariance to $\mat{I}$ and re-expresses the gradient and Hessian as
\begin{equation}
\tilde{\vect{b}} = \mat{G}^{1/2}\vect{b}, \qquad \tilde{\mat{H}} = \mat{G}^{1/2}\,\mat{H}\,\mat{G}^{1/2},
\label{eq:whitened}
\end{equation}
with $\Vlin = \tilde{\vect{b}}^{\top}\tilde{\vect{b}}$ and $\Vquad = \tfrac{1}{2}\tr(\tilde{\mat{H}}^{2})$. \emph{Rotating} by the orthogonal matrix $\mat{U}$ that diagonalises $\tilde{\mat{H}}$ then gives $\tilde{\mat{H}} = \mat{U}\mat{M}\mat{U}^{\top}$ with $\mat{M}=\mathrm{diag}(\mu_1,\ldots,\mu_n)$, and writes the gradient in the joint principal frame as $\vect{c} = \mat{U}^{\top}\tilde{\vect{b}}$. The eigenvalues $\mu_i$ are those of the matrix $\mat{H}\mat{G}$, which is the natural object here because it pairs the local curvature with the genetic covariance the population brings to it; equivalently, $\{\mu_i\}$ solve the generalised eigenvalue problem $\mat{H}\vect{v} = \mu\,\mat{G}^{-1}\vect{v}$, hence the term \emph{generalised eigenvalue rotation} for the diagonalisation. This whitening-plus-rotation construction has classical antecedents: \citet{Lande1980} used it implicitly when analysing $\mat{G}$ under stabilising selection, and \citet{MartinLenormand2006} made the generalised eigenvalue structure of $\mat{H}\mat{G}$ explicit in the context of multivariate adaptation. The contribution here is to use this rotation to expose the basis-free content of \Ag. In this rotated whitened frame the additivity index reads
\begin{equation}
\Ag = \frac{\sum_i c_i^2}{\sum_i c_i^2 + \tfrac{1}{2}\sum_i \mu_i^2}.
\label{eq:Ag_rotated}
\end{equation}
Each direction $i$ contributes $c_i^2$ to $\Vlin$ (gradient mass that the standing variance can act on) and $\tfrac{1}{2}\mu_i^2$ to $\Vquad$ (curvature mass that the standing variance feels). \Ag aggregates these direction-by-direction loads and is high when the gradient concentrates where curvature is small.

\paragraph{Scale-freeness.}
Equation~\eqref{eq:Ag_rotated} is manifestly basis-free: the $c_i$ and $\mu_i$ are intrinsic to the $(\vect{b},\mat{H},\mat{G})$ triple, not to the coordinate system used to write down those objects. The underlying invariance is that $\Vlin = \vect{b}^{\top}\mat{G}\vect{b}$ and $\Vquad = \tfrac{1}{2}\tr[(\mat{H}\mat{G})^2]$ are individually preserved under any invertible linear change of coordinates in breeding-value space (Appendix~\ref{app:variance}), so \Ag is too. The diagnostic therefore does not depend on whether traits are measured in centimetres, grams, or standardised units, nor on whether one rotates principal components before analysis. This is a desirable property for a quantity intended to be compared across studies and across systems.

\paragraph{Lines of least and greatest resistance.}
Schluter's ``genetic line of least resistance'' \citep{Schluter1996}---the leading principal axis $\vect{g}_{\max}$ of $\mat{G}$---is the direction along which evolution proceeds most readily because it carries the most variance. The rotation here adds a complementary axis: the eigenvector of $\tilde{\mat{H}}$ with the smallest $|\mu_i|$ is the direction along which the population's standing variance feels the least curvature, and therefore the direction in which prediction is most additive. When $\vect{b}$ aligns with that direction, \Ag is high; when it aligns with directions of large $|\mu_i|$, \Ag falls. A population on a fitness ridge whose $\vect{g}_{\max}$ runs along the ridge inherits both Schluter's evolvability and our additivity simultaneously; one whose variance lies across the ridge inherits neither.

% ==========================================================
\section{Selection Reshapes Genetic Covariances}
\label{sec:selection}
% ==========================================================

The additivity index depends on the genetic covariance matrix $\mat{G}$, but $\mat{G}$ is not fixed. Selection acting within a generation reshapes the distribution of breeding values among survivors, changing both the mean and the covariance. Understanding how selection transforms $\mat{G}$ is therefore essential for understanding how \Ag can change over time. We refer to the explicit within-generation mapping from $\mat{G}$ to the post-selection covariance $\Gsel$ under Gaussian stabilising selection as the \emph{Gaussian selection operator}. Coupling this selection update to the across-generation inheritance recursion yields a \emph{multi-generation loop} in which $\mat{G}$ determines \Ag, while selection and inheritance jointly reshape $\mat{G}$ (Fig.~\ref{fig:framework}).

% ----------------------------------------------------------
% TABLE 2: QUESTIONS
% ----------------------------------------------------------
\begin{table}[H]
\centering
\caption{Which quantity answers which question in the additive-channels framework.}
\label{tab:questions}
\begin{tabular}{@{}p{4.5cm}p{3.8cm}@{}}
\toprule
Question & Object(s) \\
\midrule
Is local linear prediction adequate? & $\Ag(\mat{G}, \vect{b}, \mat{H})$ \\[2pt]
Direction of short-term response? & $\Delta\bar{\vect{z}} \approx \mat{G}\vect{\beta}$ \\[2pt]
Within-generation reshaping of $\mat{G}$? & $\Gsel$ (Level 1/2 operators) \\[2pt]
Across-generation change in $\mat{G}$? & Bulmer recursion ($\rho$, $\mat{M}$) \\[2pt]
When does Gaussian closure fail? & Skewness; $\Cov(L,Q) \neq 0$ \\
\bottomrule
\end{tabular}
\end{table}


\subsection{Connecting log-fitness geometry to selection operators}

The Taylor expansion of log-fitness (Equation~\ref{eq:taylor}) can be written as:
\begin{equation}
\ell(\vect{a}) = \ell_0 + \vect{b}^\top \Delta\vect{a} + \tfrac{1}{2}\Delta\vect{a}^\top \mat{H} \Delta\vect{a}
\end{equation}
Since $W = e^\ell$, the fitness function is:
\begin{equation}
W(\vect{a}) \propto \exp\left(\vect{b}^\top \Delta\vect{a} + \tfrac{1}{2}\Delta\vect{a}^\top \mat{H} \Delta\vect{a}\right)
\label{eq:W-from-ell}
\end{equation}

For stabilising selection, the Hessian $\mat{H}$ is negative definite. Writing $\mat{H} = -\Gamma$ where $\Gamma$ is positive definite, we have:
\begin{equation}
W(\vect{a}) \propto \exp\left(\vect{b}^\top \Delta\vect{a} - \tfrac{1}{2}\Delta\vect{a}^\top \Gamma \Delta\vect{a}\right)
\label{eq:W-Gaussian}
\end{equation}

This is the form of Gaussian selection \citep{Lande1979, BartonTurelli1991}. For stabilising curvature, it is often clearer to work with $\Gamma = -\mat{H}$, a positive-definite ``curvature magnitude'' matrix; larger eigenvalues of $\Gamma$ indicate stronger selection along the corresponding directions.

We assume stabilising curvature throughout ($\Gamma \succ 0$), which ensures that the post-selection distribution is proper. When $\Gamma$ has non-positive eigenvalues---as at saddle points or under disruptive selection---the weighting in Equation~\eqref{eq:W-Gaussian} does not yield a normalisable distribution, and additional structure (truncation, finite support) would be required. This scope restriction focuses our analysis on the most common case of stabilising selection near fitness peaks, while acknowledging that other selection regimes require different mathematical treatment.

The linear term $\vect{b}^\top \Delta\vect{a}$ in Equation~\eqref{eq:W-Gaussian} shifts the mean of the post-selection distribution but does not change the covariance. This is a standard property of Gaussian updating: adding a linear term to the log-likelihood shifts the posterior mean without affecting the posterior variance \citep{Barton2017}. Therefore, we can derive covariance updates assuming centred quadratic stabilising selection ($\vect{b} = \vect{0}$) without losing generality for $\Gsel$. The additivity story involves both $\vect{b}$ and $\mat{H}$, but the selection operators need only $\Gamma = -\mat{H}$.

% ----------------------------------------------------------
% FIGURE 4: FRAMEWORK OVERVIEW
% ----------------------------------------------------------
\begin{figure*}[!htb]
\centering
\includegraphics[width=\textwidth]{figure3_framework.pdf}
\caption{Framework overview. (A)~The additive channel framework connects local geometry to multi-generation dynamics. Local geometry ($\vect{b}$, $\mat{H}$) determines both the additivity index \Ag (Eq.~\ref{eq:Ag}) and the selection operator ($\Gamma = -\mat{H}$, Eq.~\ref{eq:level1}). Selection compresses the genetic covariance $\mat{G}$, which is then transmitted through inheritance (Bulmer recursion with LD retention $\rho$ and mutation $\mat{M}$, Eqs.~\ref{eq:bulmer-genic}--\ref{eq:bulmer-LD}). The feedback loop (amber) is the key dynamic: $\mat{G}$ determines \Ag, while selection reshapes $\mat{G}$. (B)~A natural population with large genetic variance ($\tr(\mat{G}) = 3.0$) samples the curvature of the landscape, yielding $\Ag = 0.60$ (near the additive channel boundary). (C)~A breeding population with compressed variance ($\tr(\mat{G}) = 0.35$) experiences the landscape as locally flat, yielding $\Ag = 0.93$ (inside the additive channel). Natural populations and breeding programs differ in their $\rho$ and $\mat{M}$ values, leading to different equilibrium positions.}
\label{fig:framework}
\end{figure*}

\subsection{Level 1: Selection directly on breeding values}
\label{sec:level1}

Suppose breeding values follow a multivariate normal distribution before selection: $\vect{a} \sim \mathcal{N}(\vect{0}, \mat{G})$. With Gaussian fitness $W(\vect{a}) \propto \exp(-\tfrac{1}{2}\vect{a}^\top \Gamma \vect{a})$, the post-selection covariance is:
\begin{equation}
\Gsel = (\mat{G}^{-1} + \Gamma)^{-1}
\label{eq:level1}
\end{equation}
The full derivation is given in Appendix~\ref{app:selection}.

The formula has an interpretation in terms of precision (inverse covariance). The precision of the pre-selection distribution is $\mat{G}^{-1}$. Selection adds precision $\Gamma$. The post-selection precision is the sum: $(\Gsel)^{-1} = \mat{G}^{-1} + \Gamma$. Selection compresses variance by adding precision \citep{Lande1980, Bulmer1971, Bulmer1980, BartonTurelli1991}. In the isotropic case ($\mat{G} = \sigma^2 \mat{I}$, $\Gamma = \gamma \mat{I}$), this simplifies to:
\begin{equation}
\sigma^2_{\text{sel}} = \frac{\sigma^2}{1 + \gamma \sigma^2}
\label{eq:level1-scalar}
\end{equation}
Selection reduces variance, and the reduction is stronger when the initial variance is larger.

\subsection{Level 2: Selection on phenotypes}

In reality, natural selection acts on phenotypes, not directly on breeding values. If phenotype is $\vect{z} = \vect{a} + \vect{e}$, where $\vect{e}$ represents environmental deviations with $\vect{e} \sim \mathcal{N}(\vect{0}, \mat{E})$ independent of $\vect{a}$, then selection sees breeding values through environmental blur \citep{Lynch1998}. Let $\mat{P} = \mat{G} + \mat{E}$ be the phenotypic covariance, and let selection have width $\Omega$ (so $\Gamma = \Omega^{-1}$). The post-selection genetic covariance is:
\begin{equation}
\Gsel = \mat{G} - \mat{G}(\mat{P} + \Omega)^{-1}\mat{G}
\label{eq:level2}
\end{equation}

This formula shows how environmental variance buffers selection. When $\mat{E}$ is large relative to $\mat{G}$ (low heritability), the term $\mat{G}(\mat{P} + \Omega)^{-1}\mat{G}$ is small, and selection is less effective at compressing genetic variance. This is why breeding programs, which reduce environmental variance through controlled trials, can compress genetic variance more rapidly than natural selection typically does. When environmental variance vanishes ($\mat{E} \to \mat{0}$), Level 2 reduces to Level 1 with $\Gamma = \Omega^{-1}$. The Woodbury matrix identity establishes this equivalence (Appendix~\ref{app:selection}).


% ==========================================================
\section{Inheritance and Equilibrium}
\label{sec:inheritance}
% ==========================================================

Selection reshapes genetic covariances within a generation, but these reshaped covariances must be transmitted to the next generation through meiosis, recombination, and mutation. Two opposing forces govern this transmission. The first is the Bulmer effect \citep{Bulmer1971}: when selection favours specific phenotypes, it creates correlations among alleles at different loci, linkage disequilibrium (LD), that reduce additive genetic variance below what allele frequencies alone would predict. The intuition is that selection eliminates individuals with extreme breeding values, preferentially retaining those near the optimum. 

These survivors tend to carry compensatory combinations of alleles, some pushing the phenotype up and others pushing it down, thereby creating negative correlations among allele effects. The second force is recombination, which occurs in each generation and shuffles alleles during meiosis, thereby breaking up selection-induced correlations. The rate of erosion depends on recombination rates between loci: for unlinked loci, LD is halved each generation; for linked loci, erosion is slower, proportional to the recombination fraction \citep{Walsh2018}. The balance between selection that builds LD and recombination that erodes it determines the equilibrium genetic variance.

\subsection{A minimal inheritance recursion}

We decompose the genetic covariance matrix into genic variance $\Ggenic$ (depending only on allele frequencies) and LD contributions:
\begin{equation}
\mat{G} = \Ggenic + \mat{G}_{\text{LD}}
\end{equation}

A minimal recursion capturing inheritance dynamics is:
\begin{align}
\Ggenic' &= \Ggenic + \mat{M} \label{eq:bulmer-genic}\\
\mat{G}_{t+1} &= \Ggenic' + \rho(\Gsel - \Ggenic') \label{eq:bulmer-LD}
\end{align}
We refer to this across-generation update, controlled by mutational input $\mat{M}$ and LD retention $\rho$, as the \emph{inheritance recursion}. Here $\mat{M}$ is mutational input per generation, and $\rho \in [0, 1]$ is an effective LD-retention factor. When $\rho = 0$, recombination completely erases selection-induced LD. When $\rho = 1$, LD persists completely (as in clonal or selfing populations).

We emphasise the status of this recursion: it is a phenomenological closure, not a model of multilocus genetics. We retain one empirical lever ($\rho$) controlling LD persistence because the question here is how variance compression interacts with curvature, not how LD is generated in detail. The recursion captures the essential tension between selection (which creates LD and compresses variance) and recombination and mutation (which erode LD and restore variance) while remaining analytically tractable; the scope of this closure---and what a finer treatment would look like---is examined below.

\subsection{Equilibrium conditions}

At equilibrium, the genetic covariance matrix no longer changes: $\mat{G}_{t+1} = \mat{G}_t = \mat{G}^*$. This requires:
\begin{equation}
\mat{G}^* = \Ggenic^* + \rho(\mat{G}^{*,\text{sel}} - \Ggenic^*)
\label{eq:equilibrium}
\end{equation}
where $\mat{G}^{*,\text{sel}}$ is the post-selection covariance when starting from $\mat{G}^*$.

In the isotropic case with Level 1 selection, this becomes a scalar equation:
\begin{equation}
\sigma^{*2} = \sigma^2_{\text{genic}} + \rho\left(\frac{\sigma^{*2}}{1 + \gamma \sigma^{*2}} - \sigma^2_{\text{genic}}\right)
\end{equation}

Solving for the equilibrium variance $\sigma^{*2}$ yields (algebraic details are given in Appendix~\ref{app:equilibrium}):

\begin{equation}
\sigma^{*2}
= \frac{-(1-\rho)\,(1-\gamma\sigma^2_{\text{genic}}) + \sqrt{D}}{2\gamma},
\label{eq:sigma-star}
\end{equation}
where the discriminant is
\begin{equation}
D = (1-\rho)^2(1-\gamma\sigma^2_{\text{genic}})^2 + 4\gamma\,\sigma^2_{\text{genic}}(1-\rho).
\label{eq:sigma-star-discriminant}
\end{equation}

This expression, while complicated, reveals key dependencies. Equilibrium variance increases with mutational input (through $\sigma^2_{\text{genic}}$), decreases with selection strength ($\gamma$), and depends on $\rho$, typically monotonically in simple parameterisations (more LD retention leads to more persistent compression and lower equilibrium variance), though potentially non-monotone in richer models where genic variance and LD interact. Figure~\ref{fig:trajectories} shows these dynamics.

\paragraph{Beyond the scalar closure: a directional retention operator.}
A finer treatment of LD retention would replace the scalar $\rho$ with a direction-dependent retention operator $\rho_i$ that varies along the principal directions of $\mat{H}\mat{G}$ (the \textit{Generalised eigenvalue rotation and the basis-free form of \Ag} subsection). Such an operator would be generated from primitives such as recombination rates between contributing loci, the locus-by-trait genetic architecture, and the per-direction strength of selection. Computing it lies outside the scope of the additive-channels framework, which is deliberately agnostic about which loci contribute to which trait combinations. The diagnostic \Ag does not depend on $\rho$; only the equilibrium values---Equation~\eqref{eq:sigma-star} and the $\Ag^*$ derived below---do. We therefore treat the directional retention operator as a target for future work and retain the scalar closure for the analyses presented here.

\subsection{Equilibrium additivity index and the self-correcting dynamic}

At equilibrium, \Ag also reaches a stable value. Using Equation~\eqref{eq:R-isotropic}:
\begin{equation}
\Ag^* = \frac{1}{1 + n\gamma^2 \sigma^{*2}/(2|\vect{b}|^2)}
\label{eq:Ag-equilibrium}
\end{equation}
The behavior of $\Ag^*$ depends on how $\sigma^{*2}$ and $|\vect{b}|$ scale with selection strength and displacement from the optimum.

\paragraph{Scope of this equilibrium.}
The quantity $\Ag^*$ in Equation~\eqref{eq:Ag-equilibrium} is a conditional $\mat{G}$-equilibrium under a held or externally maintained local gradient $\vect{b}$ and Hessian $\mat{H}$, not a joint equilibrium of $(\bar{\vect{a}},\mat{G},\vect{b})$ on a fixed quadratic peak. The recursion that delivers $\sigma^{*2}$ holds $\vect{b}$ and $\mat{H}$ quasi-statically while $\mat{G}$ relaxes. On a fixed quadratic surface where the population is moving toward the optimum, $\vect{b}$ is itself changing each generation; the equilibrium analysis here therefore describes the regime in which $\vect{b}$ either persists (rotating-gradient or held-displacement scenarios) or in which the timescale of $\mat{G}$ relaxation is short relative to mean displacement. The case $\vect{b} = \vect{0}$, in which the population has reached the optimum, is not an equilibrium of $\Ag^*$ but its scope limit: $\Vlin \to 0$ collapses the predictive interpretation of $\Ag$ entirely, as discussed in Appendix~\ref{app:faq}.

\paragraph{The self-correcting tendency.}
During transient dynamics, selection compresses variance, which tends to raise \Ag when the gradient magnitude $|\vect{b}|$ remains substantial. The mechanism is differential scaling: $\Vquad \propto \sigma^4$ while $\Vlin \propto \sigma^2$. When selection compresses genetic variance, both components shrink, but the quadratic component shrinks \emph{faster}. A population that begins outside the additive channel (low \Ag) experiences strong curvature effects, but the very act of selection reduces variance and thereby reduces the relative importance of curvature. The population is pulled toward the channel by the same force that moves it through phenotype space.

This creates a kind of ``gravitational'' attraction toward additive regimes under directional selection on curved landscapes---not because the landscape becomes flat, but because the population's sampling of that landscape becomes confined to regions where flatness is a good approximation.

\paragraph{Transient versus equilibrium dynamics.}
The self-correcting tendency operates during transient dynamics when variance compression outpaces the decline in gradient magnitude. At equilibrium, the relationship between selection strength and \Ag is more complex. The equilibrium curvature relevance ratio $R^* = n\gamma^2\sigma^{*2}/(2|\vect{b}|^2)$ contains $\gamma^2$ in the numerator; even though stronger selection reduces $\sigma^{*2}$, the net effect depends on how steeply equilibrium variance declines with $\gamma$. Under standard mutation--selection balance, $\sigma^{*2} \propto \gamma^{-1/2}$ \citep{Lande1980, Turelli1984}, which implies $\gamma^2\sigma^{*2} \propto \gamma^{3/2}$---increasing with selection strength. The self-correcting framing therefore applies most cleanly to transient dynamics under maintained directional selection, not to equilibrium comparisons across populations with different $\gamma$.

\paragraph{When self-correction fails.}
The self-correcting tendency operates only under specific conditions. First, we require stabilising curvature ($\Gamma \succ 0$) so that selection compresses rather than inflates variance. Second, we require an approximately Gaussian distribution so that the variance decomposition remains valid. Third, and most importantly, the gradient magnitude $|\vect{b}|$ must remain substantial.

The third condition fails systematically during adaptation toward a fixed fitness optimum. For a quadratic fitness surface centered at $\vect{a}^*$, the gradient at the population mean $\bar{\vect{a}}$ is $\vect{b} = -\Gamma(\bar{\vect{a}} - \vect{a}^*)$. As selection pushes the mean toward the optimum, $|\vect{b}|$ shrinks proportionally to the displacement $\delta = |\bar{\vect{a}} - \vect{a}^*|$. In the isotropic case, the curvature relevance ratio becomes:
\begin{equation}
R = \frac{n\gamma^2\sigma^2}{2\gamma^2\delta^2} = \frac{n\sigma^2}{2\delta^2}
\label{eq:R-displacement}
\end{equation}
The selection strength $\gamma$ cancels. What matters is the ratio of variance to squared displacement. Under standard Gaussian selection, both $\sigma^2$ and $\delta^2$ shrink each generation, but $\delta^2$ shrinks \emph{faster}: within a selection step, $\delta_{\mathrm{sel}}/\delta = \sigma^2_{\mathrm{sel}}/\sigma^2 = 1/(1 + \gamma\sigma^2)$, so $\delta^2$ is scaled by the square of the same factor that scales $\sigma^2$. The ratio $\delta^2/\sigma^2$ therefore decreases monotonically and \Ag falls toward zero as the population approaches the optimum, without a preceding rise. Adaptation toward a fixed optimum produces only the late-stage fall.

The rise predicted by the self-correcting tendency requires that $|\vect{b}|$ be \emph{decoupled} from $\delta$---that is, maintained at a substantial value by some external pressure rather than collapsing as the mean approaches an optimum. Two biologically common scenarios produce this regime. First, sustained artificial selection holds $|\vect{b}|$ approximately constant for many generations because the breeder re-applies the gradient each cycle; under such selection, $\sigma^2$ compresses while $|\vect{b}|^2$ remains substantial, so $V_{\mathrm{lin}} = |\vect{b}|^2 \sigma^2$ shrinks linearly with $\sigma^2$ while $V_{\mathrm{quad}} \propto \sigma^4$ shrinks quadratically, and \Ag rises. Second, populations tracking a moving optimum---whether driven by environmental change or sustained ecological pressure---experience an approximately constant gradient as long as their lag behind the optimum stays roughly constant. In both cases, when the maintaining pressure eventually relaxes, $|\vect{b}| \to 0$ and \Ag falls. Figure~\ref{fig:adaptation-trajectory} illustrates this rise-then-fall under maintained directional selection.

The crossover between the rising and falling phases occurs when $|\vect{b}|$ decays to a threshold below which curvature once again dominates. For \Ag to exceed a threshold $\alpha$, we require $R < (1-\alpha)/\alpha$, which in the maintained-selection regime translates to a condition on the gradient magnitude:
\begin{equation}
|\vect{b}|^2 > \frac{n\gamma^2 \sigma^2}{2}\,\frac{\alpha}{1-\alpha}
\label{eq:crossover-displacement}
\end{equation}
While the maintaining pressure keeps $|\vect{b}|$ above this threshold, the population stays in the additive channel; once it falls below, \Ag drops. The duration of the channel therefore depends on the timescale of selection-pressure decay relative to the timescale of variance compression.

% ----------------------------------------------------------
% FIGURE: ADAPTATION TRAJECTORY
% ----------------------------------------------------------
\begin{figure*}[!htb]
\centering
\includegraphics[width=\textwidth]{figure4_trajectory.pdf}
\caption{\textbf{Additivity index dynamics under maintained directional selection that eventually relaxes.} A population experiences a sustained directional gradient $|\vect{b}|$ that decays slowly over time, while stabilising curvature $\gamma$ compresses genetic variance $\sigma^2$ on a faster timescale. Dynamics computed from the Bulmer-style inheritance recursion of Eqs.~\ref{eq:bulmer-genic}--\ref{eq:bulmer-LD} coupled to Gaussian selection (Eq.~\ref{eq:level1}); $\Ag$ is evaluated at each generation via Eq.~\ref{eq:Ag}. This setup represents either sustained artificial selection in a breeding programme or response to an environmental pressure that gradually subsides. \textbf{(A)}~Genetic standard deviation $\sigma$ (solid green) and gradient magnitude $|\vect{b}|$ (dashed orange) over generations, log time axis. Variance compresses to its equilibrium within $\sim 10$ generations; the gradient decays an order of magnitude slower. \textbf{(B)}~Additivity index \Ag rises from 0.47 to a peak of 0.78 at generation 10 (filled circle) as variance compresses while the gradient remains substantial. The population enters the additive channel ($\Ag \geq 0.75$, shaded) at generation 4 and remains there for 26 generations, then falls toward zero as $|\vect{b}| \to 0$. \textbf{(C)}~Phase portrait in $(|\vect{b}|, \sigma)$ space with \Ag as a colour field. The trajectory (orange) moves vertically downward as variance compresses, plateaus near the additive-channel boundary, then sweeps left as the gradient decays, ending in the low-\Ag region near $|\vect{b}| = 0$. The same algebra runs in reverse for adaptation toward a fixed optimum (where $\vect{b} = -\Gamma\bm{\delta}$ collapses with $\bm{\delta}$): only the late-stage fall is observed, without a preceding rise. Parameters: $n = 4$, $\gamma = 0.05$, $\rho = 0.95$, $M = 10^{-3}$, $\sigma^2_0 = 20$, $|\vect{b}|_0 = 0.3$, gradient decay rate $\lambda_b = 0.005$ per generation.}
\label{fig:adaptation-trajectory}
\end{figure*}

\paragraph{Biological interpretation.}
The failure of self-correction near optima is not a problem for the framework but a clarification of its scope. The additivity index diagnoses when linear prediction of directional response is accurate. Near a fitness optimum where $\vect{b} \approx 0$, there is no directional response to predict: both the additive prediction $\Delta\bar{\vect{z}} = \vect{G}\bm{\beta}$ and the true response are near zero. Low \Ag at equilibrium reflects the absence of directional selection, not a failure of additive genetics.

The practical implication is that \Ag is most informative during active adaptation: artificial selection in breeding programs, response to environmental change, range expansion, or any scenario where populations track a moving or displaced optimum. The framework's prediction that maintained directional selection sustains the additive channel for many generations is consistent with the Illinois long-term selection experiment for oil and protein content in maize, where additive response has continued for over a hundred generations of artificial selection \citep{Dudley2007}, a regime where the gradient is re-applied each cycle by the breeder rather than collapsing toward an internal optimum. For populations at mutation--selection balance near a stable optimum, the relevant question shifts from prediction accuracy to variance maintenance and stability---questions addressed by classical theory \citep{Lande1980, Turelli1984, Burger2000} rather than by the additivity index.

\paragraph{Contraction versus expansion.}
In this local-Gaussian setting, additive channels reflect a balance between processes that \emph{contract} the neighbourhood sampled by a population (stabilising selection and selection-generated LD, aided by high LD retention) and processes that \emph{expand} or restore that neighbourhood (mutation, migration, and recombination), with drift typically reducing genic variance and thereby also favouring contraction. When contraction dominates, variance shrinks and the population experiences the landscape as locally flat; when expansion dominates, the variance cloud ``tastes'' the curvature and \Ag falls. This framing parallels standard mutation--selection balance but shifts the question from ``how much variance is maintained?''\ to ``what is the geometry of that variance relative to the fitness surface?'' Elite breeding programs operate deep in the additive channel because of inbreeding (high $\rho$) and strong artificial selection; natural outcrossing populations occupy an intermediate zone; wide crosses can temporarily push populations toward lower \Ag until selection re-compresses variance. A heuristic regime map summarising these contrasts is provided in the SI.

% ----------------------------------------------------------
% TABLE tab:Ag_interpretation: REMOVED in this revision.
% Its three-row threshold content is already covered in prose; the table
% reproduced material with no added clarity.
% ----------------------------------------------------------

\subsection{Natural populations versus breeding programs}

Figure~\ref{fig:trajectories} shows how variance and \Ag evolve under different parameter regimes. In natural populations, mutation continuously restores variance ($\mat{M} > 0$) and outcrossing erodes LD ($\rho$ moderate), so $\mat{G}$ equilibrates near the mutation--drift balance $N_e\mat{M}$ \citep{Turelli1988, TurelliBarton1994}. The trajectory of \Ag, however, depends on the temporal structure of the directional gradient $\vect{b}$. Under sustained directional pressure, \Ag settles at a moderate equilibrium value bounded away from one because $\Vquad \propto \|\mat{G}\|^2$ does not collapse. Under the episodic pressure that natural populations more typically experience---bouts of climatic, demographic, or biotic selection separated by stretches of mutation--drift balance---\Ag spikes upward during each bout and falls back near zero between them. In breeding programs, by contrast, mutational input is negligible on breeding timescales ($\mat{M} \approx 0$), inbreeding and selfing increase $\rho$, and directional selection is sustained across cycles. The result is progressive variance collapse and monotonically increasing \Ag: breeding populations contract into the additive channel and remain there, often for hundreds of generations, while natural populations only briefly visit it \citep{Technow2021}.

% ----------------------------------------------------------
% FIGURE 5: NATURAL VS BREEDING
% ----------------------------------------------------------
\begin{figure*}[!htb]
\centering
\includegraphics[width=\textwidth]{figure5_natural_vs_breeding.pdf}
% NEW CAPTION (sustained-vs-episodic, with specific values matching the rendered figure):
\caption{\textbf{Trajectories under sustained versus episodic directional selection.} Both populations evolve from identical initial conditions ($\mat{G}_0 = 2\mat{I}$, $\Ag(0) = 0.50$) under the inheritance recursion of the \textit{Inheritance and Equilibrium} section. The breeding population (vermillion) experiences \emph{sustained} directional pressure ($\vect{b}$ constant). The natural population (blue) experiences three log-spaced episodes of directional pressure (centres at generations~40, 240, 700; light-blue bands) with weak baseline pressure between bouts. \textbf{(A)}~Mean genetic variance per axis on log--log axes. Breeding's $\mat{G}$ collapses (Bulmer compression with negligible $\mat{M}$); natural's $\mat{G}$ relaxes monotonically toward its $N_e\mat{M}$ equilibrium. \textbf{(B)}~Additivity index. Breeding crosses the illustrative channel threshold ($\Ag \geq 0.8$, shaded yellow) at generation~40 and remains inside through the end of the run (gen~1000). Natural's \Ag spikes during each episode (peaks $0.53$, $0.59$, $0.65$) but falls to $\approx 0$ between bouts, where weak baseline $\vect{b}$ makes $\Vlin \ll \Vquad$. \textbf{(C)}~Natural variance decomposition: $\Vlin = \vect{b}^\top\mat{G}\vect{b}$ (Eq.~\ref{eq:Vlin}) surges with $\|\vect{b}\|$ during episodes, while $\Vquad = \tfrac{1}{2}\tr[(\mat{H}\mat{G})^2]$ (Eq.~\ref{eq:Vquad}) depends only on $\mat{G}$ and declines smoothly. \textbf{(D)}~Breeding variance decomposition on log~$y$: $\Vquad \sim \|\mat{G}\|^2$ collapses faster than $\Vlin \sim \|\mat{G}\|$---the mechanism behind $\Ag \to 1$. The dotted vertical at gen~100 marks the Illinois long-term selection experiment horizon. Parameters: $n = 4$, $\mat{\Gamma} = 0.20\mat{I}$, $\vect{b}_{\text{full}} = (0.30, 0.20, 0.15, 0.10)$; breeding ($\mat{E} = 0.05\mat{I}$, $\mat{M} = 10^{-5}\mat{I}$, $\rho = 0.95$, $N_e = 100$); natural ($\mat{E} = 0.50\mat{I}$, $\mat{M} = 2{\times}10^{-3}\mat{I}$, $\rho = 0.40$, $N_e = 500$; episode envelope $0.05 \to 1.0$). The natural-population update uses the Lynch--Hill mutation--drift closure $\mat{G}_{\text{genic}}' = \mat{G}_{\text{genic}} + \mat{M} - \mat{G}_{\text{genic}}/N_e$ to give a well-defined long-term equilibrium.}
%
% OLD CAPTION (kept here for comparison, not active):
% \caption{Natural versus breeding dynamics. Contrasting evolutionary trajectories under different parameter regimes. (A)~Genetic variance over generations: natural populations equilibrate at moderate variance due to mutation-selection balance; breeding populations collapse toward zero due to negligible mutation and high LD retention. (B)~Additivity index trajectories: breeding programs enter the additive channel ($\Ag > 0.8$, shaded as illustrative threshold) while natural populations stabilise at moderate \Ag. (C,D)~Variance decomposition into $\Vlin$ and $\Vquad$ components. In breeding programs, $\Vquad \to 0$ as variance collapses, making curvature irrelevant for the local neighbourhood sampled. Parameters shown are illustrative toy values: natural ($M = 0.02$, $\rho = 0.4$); breeding ($M = 0.001$, $\rho = 0.85$).}
\label{fig:trajectories}
\end{figure*}

% ----------------------------------------------------------
% TABLE tab:regimes: REMOVED in this revision.
% Its row-by-row regime contrasts are now carried by fig:trajectories and
% the surrounding natural-vs-breeding prose. Moved to SI if needed.
% ----------------------------------------------------------

\paragraph{Scope of the equilibrium analysis.}
The equilibrium expressions derived above (Equations~\eqref{eq:sigma-star} and \eqref{eq:Ag-equilibrium}) describe where populations settle under long-term mutation--selection balance. This analysis complements, rather than replaces, the transient dynamics that are the framework's primary focus. For populations actively adapting---whether through artificial selection, response to environmental change, or range expansion---the transient behavior of $\mathcal{A}_g$ is the relevant diagnostic. The equilibrium analysis becomes relevant when asking why natural populations maintain the variance structures they do, and when comparing populations that have reached approximate stationarity under different selective regimes. For populations near equilibrium under stabilising selection, the additivity index may be low (because $|\bm{b}| \approx 0$), but this reflects the absence of directional response to predict rather than a breakdown of additive genetics.
% ==========================================================
\section{Dimensional Considerations}
\label{sec:dimension}
% ==========================================================

The curvature relevance ratio $R$ (Equation~\ref{eq:R-isotropic}) contains a factor of $n$, the number of dimensions, raising the question of whether high-dimensional systems behave fundamentally differently from low-dimensional ones. Our operational conclusion is that the relevant dimension is the dimension actually explored by $\mat{G}$ near the current mean, which is typically far smaller than the nominal genotype space.

At first glance, high dimensionality appears problematic: with all parameters except $n$ held fixed, $R \propto n$ and \Ag declines. This naive scaling assumes that $|\vect{b}|$, $\gamma$, and $\sigma^2$ are independent of $n$, which biological systems violate. Consider a population displaced from an optimum by $\delta$ in each of the $n$ dimensions. The gradient magnitude scales as $|\vect{b}| \sim \sqrt{n}\gamma\delta$, since the gradient is proportional to displacement times curvature, summed in quadrature across dimensions. Substituting,
\begin{equation}
R = \frac{n\gamma^2\sigma^2}{2 n\gamma^2\delta^2} = \frac{\sigma^2}{2\delta^2},
\end{equation}
and dimension cancels.

This cancellation is a special case, not a general theorem. It rests on isotropic variance, isotropic curvature, and uniform displacement across all dimensions. Two departures break it. First, if the population is displaced from the optimum in only $k < n$ dimensions (concentrated displacement), $|\vect{b}|^2 = k\gamma^2\delta^2$ and $R = n\sigma^2/(2k\delta^2)$: the ratio $n/k$ reappears. Second, if curvature is concentrated along directions of high genetic variance, $\Vquad$ is amplified relative to the isotropic case. The general principle is that three effective dimensionalities matter---directions carrying substantial variance, curvature, and gradient---and that their alignment, not the nominal dimension, governs $R$. When variance lies along directions of weak curvature while the gradient lies along directions of strong curvature, dimensional structure can either ameliorate or exacerbate curvature effects.

A natural way to quantify the variance side is the participation ratio,
\begin{equation}
\keff = \frac{(\tr \mat{G})^2}{\tr(\mat{G}^2)} = \frac{\left(\sum_i \lambda_i\right)^2}{\sum_i \lambda_i^2},
\label{eq:keff}
\end{equation}
where $\lambda_i$ are the eigenvalues of $\mat{G}$. For a matrix with $k$ equal nonzero eigenvalues, $\keff = k$; for exponentially decaying spectra, $\keff \ll n$. Real $\mat{G}$ matrices are typically dominated by a few principal components \citep{Blows2007, Kirkpatrick2009, Walsh2009}, so the population samples only a low-dimensional subspace of the fitness surface and curvature in unused dimensions does not contribute to $\Vquad$. Replacing $n$ with $\keff$ gives an effective curvature relevance ratio,
\begin{equation}
R \approx \frac{\keff\, \gamma^2 \sigma^2}{2 |\vect{b}|^2},
\label{eq:R-effective}
\end{equation}
in which the relevant curvature is that of the subspace the population actually explores. The naive $n$-scaling is therefore a cautionary strawman that applies only under the implausible assumption of isotropic variance in a high-dimensional space. The arguments here are heuristic; a rigorous treatment of dimensional behaviour in realistic high-dimensional systems remains open.

\paragraph{On the asymmetric anisotropy in this argument.}
Equations~\eqref{eq:R-isotropic} and~\eqref{eq:R-effective} treat $\mat{G}$ as anisotropic (via $\keff$) while keeping $\mat{H}=-\gamma\mat{I}$ isotropic. This asymmetry is didactic: it isolates the dimensional structure of $\mat{G}$ from the dimensional structure of $\mat{H}$, so the participation-ratio argument shows transparently that the variance side, not the curvature side, sets the operative dimension. The fully anisotropic case is treated formally in the \textit{Generalised eigenvalue rotation and the basis-free form of \Ag} subsection, where the simultaneous diagonalisation of $(\mat{G},\mat{H})$ via the generalised eigenvalue problem $\mat{H}\vect{v}=\mu\,\mat{G}^{-1}\vect{v}$ gives $\Vquad = \tfrac{1}{2}\sum_i \mu_i^2$ in terms of the joint eigenvalues. In that frame $\Vquad$ depends on the alignment of the principal directions of $\mat{H}$ and $\mat{G}$, not on either spectrum alone, and the participation-ratio formula above is a special case in which curvature is direction-uniform. Empirically, anisotropic curvature can arise from gene network architecture and emerges naturally in simulations that resolve $G$-to-$P$ structure: \citet{Technow2026} report, in $E(NK)$ simulations of breeding cycles, an effective curvature whose principal axes shift across replicate selection trajectories despite identical population structure, an instance of what an anisotropic and history-dependent $\mat{H}$ would look like in operation.

% ==========================================================
\section{Discussion}
\label{sec:discussion}
% ==========================================================

We have developed a framework for thinking about when additive genetic models provide adequate approximations to population movement on a fitness landscape. The central concept is the additive channel: a region of genetic space in which the population's spread interacts with the log-fitness landscape, so that curvature contributes little to fitness variance. When $\Vquad$ is small relative to $\Vlin$, linear ranking and local linear response approximations should be accurate, and additive models, whether for genomic prediction, breeding value estimation, or short-term evolutionary forecasting, should perform well. 

In this view, additive models are reasonably interpreted as tangent-plane descriptions, an interpretation developed explicitly by \citet{Rice2008} and \citet{MiloccoSalazarCiudad2022} for the case of nonlinear genotype--phenotype maps: they can be accurate when the population samples a small region of phenotype space, but they degrade when evolution or breeding forces the population to explore a sufficiently large region that slope and curvature vary across the sampled area. The contribution of $\Ag$ here is to make the locality condition quantitative.

This framework offers a geometric language for discussing additivity. Rather than asking whether genes ``act additively,'' we ask whether the population's genetic variance lies in directions where the log-fitness surface is relatively flat. This shift in perspective may help reconcile molecular evidence for widespread gene interactions with statistical evidence that additive models often suffice \citep{HillGoddardVisscher2008, SellaBarton2019}.

The framework connects local geometry to multi-generation dynamics through a unified treatment. The log-fitness Taylor expansion provides both the additivity diagnostic (via $\vect{b}$ and $\mat{H}$) and the selection operators (via $\Gamma = -\mat{H}$). This approach builds on classical results on Gaussian selection \citep{Lande1979, BartonTurelli1991, Barton2017} and provides a diagnostic lens for assessing when those results apply. The framework also suggests mechanisms for additive channel persistence and exit. Selection compresses variance, which tends to raise \Ag under the conditions we have described, creating partial self-correction. This dynamic connects to work on the Bulmer effect \citep{Bulmer1971} and variance maintenance \citep{Turelli1988, Burger2000}.

\subsection{Connections to breeding programs}

The framework may help explain why breeding programs achieve reasonable prediction accuracy despite underlying epistasis \citep{Dudley2007, Dudley2009}. Breeding compresses variance through repeated selection and increases LD retention through inbreeding, driving populations toward high-\Ag regimes \citep{Technow2021, CooperMessina2021}. This perspective aligns with observations from long-term selection experiments. The Illinois long-term selection experiment for oil and protein content in maize has continued responding to selection for over 100 generations \citep{Dudley2007}. The framework suggests this may occur because variance compression keeps populations in locally additive regimes even as they traverse globally curved landscapes. 

Cooper first observed that the rise of \%GCA under $K = 15$ isolation in the $E(NK)$ simulations of \citet{Technow2021} resembled the variance dynamics of our framework. The geometric reading developed here makes that resemblance precise: the conversion of variance under drift is, locally, the contraction of $\mat{G}$ along the principal directions of $\mat{H}\mat{G}$, with the rate of conversion in any given direction set by how much variance retention that direction enjoys.

Modern genomic selection methods \citep{Meuwissen2001, VanRaden2008} implicitly rely on additive models, yet achieve substantial prediction accuracy. The additive channel concept suggests why: breeding populations contract into high-\Ag regimes where additive approximations capture most fitness-relevant variation. The challenge of introducing genetic diversity from unadapted germplasm \citep{Jordan2011, Sanchez2023} may partly reflect difficulty maintaining high \Ag when variance is suddenly expanded. 

Work on crop improvement under climate change \citep{Messina2023, CooperMessina2021} emphasises the importance of understanding how genetic architecture affects response to selection. The additive channel framework may provide a conceptual tool for thinking about when linear prediction will work well and when it may break down; for instance, when introgressing exotic germplasm or when selection intensities change.

A related implication is that a breeding program may plateau even while remaining deep in an additive channel. High \Ag indicates that local curvature contributes little to fitness variance in the currently sampled neighbourhood so that additive prediction can remain reliable; it does not, by itself, guarantee sustained response. In our notation, short-term response is governed by the local slope and the genetic variance aligned with it, $\Vlin=\vect{b}^\top\mat{G}\vect{b}$. Repeated selection can shrink $\mat{G}$ (and generate LD that depresses observable additive variance), while continued improvement can also reduce the gradient magnitude as the population approaches a local optimum or ridge, $\|\vect{b}\|\to 0$. Either process drives $\Vlin$ downward and can stall gains even when \Ag is high. This idea clarifies a common breeding experience: introgression or wide crosses can restart long-run progress by re-expanding variance and introducing new combinations, but they may transiently reduce \Ag and prediction accuracy because the population begins sampling a broader, more curved section of the landscape; subsequent cycles then re-compress variance and can restore high-\Ag, high-predictability behavior.

% ----------------------------------------------------------
% ----------------------------------------------------------
% FIGURE fig:keff: REMOVED in this revision.
% The tightened Dimensional Considerations section carries the participation-
% ratio point in prose without a supporting figure. Moved to SI if needed.
% ----------------------------------------------------------

\subsection*{Model selection in genomic prediction}

The additivity index suggests a practical criterion for deciding when epistatic models warrant investment. From the definition $\mathcal{A}_g = V_{\text{lin}}/(V_{\text{lin}} + V_{\text{quad}})$, the fraction of genetic log-fitness variance missed by an additive model is exactly $1 - \mathcal{A}_g$. This quantity bounds the maximum improvement available from adding epistatic terms.

Consider a breeder deciding whether to move from additive GBLUP to a model that captures epistasis, such as reproducing kernel Hilbert space (RKHS) regression or explicit dominance and epistatic terms. The potential gain in explained variance is at most $V_{\text{quad}}$. But this theoretical gain must be weighed against estimation cost: epistatic models require estimating $O(p^2)$ parameters from the same data, where $p$ is the number of markers. Unless the training population is large relative to this parameter count, gains in model fidelity may be offset by increased estimation error.

A rough decision heuristic follows from this logic. When $\mathcal{A}_g > 0.95$, the maximum gain from epistatic terms is less than 5\% of log-fitness variance; additive models are almost certainly sufficient, and the estimation cost of epistatic models is unlikely to be justified. When $\mathcal{A}_g$ falls between 0.85 and 0.95, additive models remain adequate for most purposes, though epistatic models may provide marginal improvement in very large training populations. When $\mathcal{A}_g$ falls below 0.85, curvature contributes meaningfully to fitness variation, and epistatic models may improve prediction---provided training data are sufficient to estimate the additional parameters reliably.

This framing clarifies why additive GBLUP succeeds in elite breeding programs despite molecular evidence for gene interactions. Variance compression pushes these populations into high-$\mathcal{A}_g$ regimes where the theoretical ceiling for epistatic improvement is low. The same framing suggests that wide crosses or introgression from unadapted germplasm may temporarily reduce $\mathcal{A}_g$ by expanding the variance cloud into regions where curvature matters, potentially justifying investment in epistatic models during early cycles before selection re-compresses variance.

We emphasise that $\mathcal{A}_g$ bounds theoretical improvement, not realised improvement. Estimation error, model misspecification, and overfitting all reduce actual gains below this bound. The appropriate threshold for investing in epistatic models therefore depends on training population size, marker density, and computational resources, not on $\mathcal{A}_g$ alone. Nevertheless, a population with $\mathcal{A}_g > 0.95$ offers little room for improvement regardless of these factors, while one with $\mathcal{A}_g < 0.7$ presents a substantial target that sophisticated models might partially capture.
\subsection{Connections to natural populations}

In natural populations the picture differs. Mutation continuously restores variance, and outcrossing erodes selection-induced LD; the balance between these forces determines equilibrium \Ag, which may stabilise at moderate values \citep{Turelli1988, TurelliBarton1994}. This reading is consistent with the resolution of the ``missing heritability'' puzzle \citep{Manolio2009} through common variants of small effect \citep{Yang2010, Yang2015, LopezCortegano2019}: if outcrossing populations sit in moderate-\Ag regimes, additive variance will dominate observed responses even when the genotype--phenotype map involves abundant interactions.

Four concrete uses for evolutionary biologists follow from this picture. First, \Ag converts the qualifiers under which Lande's equation, Schluter's lines of least resistance, and Fisher's fundamental theorem are known to operate into a quantity computable from the same $(\vect{b},\mat{H},\mat{G})$ that empirical studies already attempt to estimate. A selection-gradient study reporting $\bm{\beta}$ and $\bm{\gamma}$ has, in our notation, the ingredients for $\Vlin$ and $\Vquad$; reporting \Ag tells the reader how much of the genetic log-fitness variance the linear analysis actually captures. Second, \Ag is a direct partner to the Hansen--Houle evolvability metric $e(\vect{b}) = \vect{b}^\top \mat{G} \vect{b}$ \citep{HansenHoule2008}, since $e(\vect{b}) = \Vlin$ in our notation. Evolvability quantifies the magnitude of additive variance in the gradient direction; \Ag quantifies whether that variance will actually govern response, or whether curvature will displace it. Third, the framework predicts that populations near a fitness optimum are \emph{less} reliably additive than populations far from one. Because $\Vlin$ is quadratic in $\vect{b}$ while $\Vquad$ does not depend on $\vect{b}$, the stable populations that dominate the empirical selection-gradient literature are precisely those in which $\vect{b}\to\vect{0}$ collapses $\Vlin$ without collapsing $\Vquad$. Low \Ag here does not mean additive genetics has broken down---it means there is little directional variance for additive analysis to capture, and the corresponding cautions apply to inferences drawn from such studies. Fourth, \Ag provides a comparative tool. Adaptive radiations are predicted to be episodes of high \Ag along $\vect{g}_{\max}$ during the directional phase of niche occupation, falling as lineages settle near new optima---complementing Schluter's account of \emph{which} directions adaptive radiations follow \citep{Schluter1996} with an account of \emph{when} they behave additively along those directions. The same frame applies to populations tracking rotating gradients under environmental change, where alignment between $\vect{g}_{\max}$ and the new $\vect{b}$ becomes the key determinant of additive predictability.

We restrict this rotating-$\vect{b}$ interpretation to environmental shifts that re-weight existing selective pressures. Some shifts do not rotate the existing fitness landscape but enlarge it---a novel pathogen, an unfamiliar climate, a new herbicide---opening fitness-relevant directions along which the population carries no standing variance. The local-geometric framework as currently formulated does not address this case.

A simulation-based validation of the central identity $\Ag = R^2$, including its behaviour across stochastic replicates and the across-replicate reproducibility of the geometric quantities on which the framework depends, is reported in Appendix~\ref{app:simulation}.

\subsection{Mechanistic versus statistical epistasis}
\label{sec:mech_vs_stat}

A subtlety threads the framework and should be made explicit. Two senses of ``epistasis'' coexist in the literature, and the additive-channels diagnostic measures one while leaving the other untouched. \emph{Mechanistic} epistasis is a property of the genotype-to-phenotype map: the effect of one allele on the phenotype depends on the genotype at another locus. This is basic biology, including regulatory networks, metabolic pathways, and protein--protein interactions, and it does not disappear merely because a population is well described by an additive statistical model. The other sense is \emph{statistical} epistasis, which is a property of the variance partition in a particular population at a particular allele-frequency configuration: the fraction of among-individual variance not captured by an additive model. The two are distinct, and our framework's prediction that elite breeding populations operate deep in the additive channel does not require, or imply, that mechanistic epistasis has somehow disappeared from those populations. Recent extensions of this line of work \citep{Morrissey2015, GonzalezForero2026} further develop the consequences of $G\!\to\!P$-map nonlinearity for selection-response prediction beyond what we treat here; the additive-channels diagnostic operates downstream of any such mapping, on the breeding-value cloud it ultimately delivers.

The bridge between the two is straightforward and worth stating cleanly. A locus contributes to the variance partition only if it is segregating in the population being evaluated. An interaction between two loci contributes to non-additive variance only if \emph{both} partners are segregating; a regulatory interaction whose other arm is fixed---by drift, by purifying selection, or by long-term conservation---generates no variance among individuals, even though the molecular interaction is biologically active in every cell of every individual. Mechanistic epistasis lives upstream in the genotype-to-phenotype map, where allelic effects depend on genetic background. In the framework presented here, the fitness-relevant nonlinear geometry of the BV-to-fitness map is captured by the local Hessian $\mat{H}$, and the part of that geometry visible to the standing breeding-value cloud is governed by $\mat{H}\mat{G}$ through $\Vquad$. When alleles at one partner of a mechanistic interaction are fixed, $\mat{G}$ loses rank along that direction, the corresponding eigenvalues of $\mat{H}\mat{G}$ collapse, and $\Vquad = \tfrac{1}{2}\tr[(\mat{H}\mat{G})^2]$ shrinks accordingly. The interaction has been silenced from the variance partition without being silenced from the biology.

Several mechanisms drive populations toward this silencing in different parts of trait space. Drift fixes alleles in finite populations, particularly those isolated by mating-system barriers (e.g., heterotic groups in maize breeding); the segregating spectrum thins each generation, and the surviving variance increasingly comes from interactions whose partners have been removed from circulation \citep{Technow2021}. Purifying selection at conserved core regulators (e.g., transcription factors essential to angiosperm development for hundreds of millions of years) silences the same interactions across an entire clade by fixing the partner. Sustained directional selection compresses $\mat{G}$ along high-curvature directions, removing variance from regions where the population would otherwise feel non-additive structure. All three mechanisms produce the same geometric outcome: a thinner, more directionally-constrained $\mat{G}$, paired with the same $\mat{H}$, yielding a lower $\Vquad$ and a higher \Ag.

\subsubsection*{An intuitive bridge to the GCA-fraction}

This connection has a familiar empirical face. In hybrid breeding programs, the standard variance partition decomposes hybrid trait variance into general combining ability (GCA, the additive contribution attributable separately to each parental line) and specific combining ability (SCA, the residual from line-by-line interactions). The breeder's diagnostic is the GCA-fraction $\eta = (V_{\mathrm{GCA}_1} + V_{\mathrm{GCA}_2})/V_G$, which empirical work has long shown rises as drift fixes alleles within heterotic groups \citep{Technow2021}. We propose that $\eta$ is a multilocus statistical shadow of \Ag, expressed in two different languages, rather than its formal equivalent. The intuition runs as follows. GCA captures, by construction, the variance attributable to the additive substitution effects of alleles transmitted from each heterotic group, projected onto the hybrid population, the discrete, finite-locus analogue of $\Vlin = \vect{b}^\top \mat{G} \vect{b}$, which captures the variance attributable to the linear gradient of log-fitness in breeding-value space. Symmetrically, SCA captures the residual variance from interactions between parental contributions and across loci, the analogue of $\Vquad = \tfrac{1}{2}\tr[(\mat{H}\mat{G})^2]$, which captures variance attributable to the local curvature of the log-fitness surface. The framework's claim that $\Vquad$ shrinks faster than $\Vlin$ as $\mat{G}$ contracts maps directly onto the empirical observation that SCA collapses faster than GCA as drift thins the segregating spectrum, because both partners of an interaction must be variable for the interaction to contribute to among-individual variance. We expect $\eta$ and \Ag to converge quantitatively under three conditions: (i)~breeding-value distributions are approximately Gaussian, so the variance identity $\Cov(L,Q)=0$ holds; (ii)~the second-order Taylor expansion of phenotype around the population mean breeding value adequately captures the non-additive structure, with cubic-and-higher terms small; and (iii)~the two heterotic groups contribute symmetrically, so the hybrid distribution inherits the assumed structure cleanly. Each violation of these conditions creates a specific, identifiable gap between the empirical $\eta$ and the analytical \Ag, and quantifying these gaps---through both rigorous derivation and simulation in the full heterotic-group setup---is a natural extension we leave to future work. We present the connection in this paper as a structural correspondence and a guide to empirical interpretation, not as a claim of formal equivalence.


\subsubsection*{Relation to Lande (1980): mutation--selection balance and the structure of $\mat{G}$}

A natural point of comparison is \citet{Lande1980}, who asked why additive genetic variances and covariances persist under stabilising selection and what determines their structure. In Lande's model, pleiotropic mutations introduce a mutational covariance matrix (often denoted $\mat{M}$), while stabilising selection removes variation according to the local curvature of the fitness surface (equivalently $\Gamma=-\mat{H}$ in our notation). Under Gaussian approximations and in regimes where linkage disequilibrium is negligible in large, approximately random-mating populations, the equilibrium $\mat{G}$ reflects a balance between mutational input and the curvature of selection, with more variance accumulating in directions that are relatively weakly penalised by curvature \citep{Lande1980}. This provides a mechanistic explanation for why $\mat{G}$ has a characteristic shape in natural populations even when selection is stabilising.

Our framework is complementary rather than a re-derivation. Lande (1980) primarily concerns the \emph{maintenance and equilibrium structure} of $\mat{G}$ under stabilising selection and mutation. Here we focus on a different question: given the current $(\mat{G},\vect{b},\mat{H})$ at a particular point along an evolutionary trajectory, is the local linear approximation to log-fitness adequate for ranking and prediction? This is precisely what the additivity index $\Ag$ quantifies. Notably, the diagnostic is most informative in off-optimum settings with $\vect{b}\neq\vect{0}$. In contrast, classic stabilising-selection equilibria are often analysed at or near an optimum where $\vect{b}\approx\vect{0}$ and curvature dominates local change.

Finally, the relationship between Lande's ``LD-negligible'' regime and our inheritance closure is direct. Lande (1980) notes that when recombination dominates selection on polygenic variation, linkage disequilibrium is expected to be small and pleiotropy can be the main source of genetic covariances \citep{Lande1980}. In our notation, this corresponds to small effective LD retention (low $\rho$), under which $\mat{G}$ remains close to its genic component and mutation--selection balance provides a natural baseline for the magnitude and orientation of $\mat{G}$. Conversely, in breeding programs (high $\rho$ and negligible $\mat{M}$ on the breeding timescale), selection-induced variance compression can push populations into high-$\Ag$ regimes even far from long-term mutation--selection equilibrium. This contrast helps explain why additive prediction can be successful in both natural and applied contexts, but for different mechanistic reasons.


\subsection{What the framework does not offer}

We have flagged the main scope limitations in the introduction; two deserve further emphasis here. First, the framework is silent about whether any particular system will have high or low \Ag: the actual value depends on the specific landscape and population, and must be measured or estimated empirically. Second, when breeding-value distributions deviate substantially from Gaussianity---under strong selection on few loci, multimodal distributions, or genuinely heavy-tailed architectures---the Gaussian closures we use lose their exact form, even though \Ag remains a monotone diagnostic across the realistic genetic architectures we tested (Figure~\ref{fig:Ag_validation}). The dimensional analysis is heuristic rather than rigorous, and understanding how additive channels behave in high-dimensional systems with realistic genetic architectures remains open.

\subsection{Empirical challenges and a measurement pathway}

Estimating \Ag empirically requires measuring $\mat{G}$, the gradient $\vect{b}$, and the Hessian $\mat{H}$. The first two are routinely estimated in quantitative genetics: $\mat{G}$ from pedigree or genomic data, and $\vect{b}$ from selection differentials or from breeding values regressed on fitness \citep{LandeArnold1983}. The Hessian is harder, since it requires second derivatives of fitness, but quadratic regression of fitness on phenotypes or breeding values gives the necessary estimates \citep{PhillipsArnold1989, BlowsBrooks2003}. Recent advances in high-throughput genotype--fitness mapping---deep mutational scanning \citep{Fowler2014}, combinatorial mutagenesis \citep{Starr2020}, and CRISPR-based screens---now provide $\vect{b}$ and $\mat{H}$ at unprecedented resolution for proteins, regulatory elements, and small genome regions; combined with population-genetic estimates of $\mat{G}$, they offer a concrete pathway for computing \Ag in real systems. In agricultural contexts, breeding populations with detailed pedigree records and multi-environment yield trials offer decades of data from which all three ingredients can be extracted \citep{CooperMessina2014}, and experimental evolution in microbes \citep{Lenski2017} provides another tractable testbed. A practical workflow is to compute $\Vlin$ from estimated $\mat{G}$ and $\vect{b}$, approximate $\Vquad$ from a fitted local quadratic term, and ask whether prediction accuracy tracks the resulting \Ag across time, environments, or breeding cycles.

\subsection{Open questions}

Several questions remain for future work. How do additive channels behave in high-dimensional systems with realistic genetic architectures? What is the typical \Ag of natural populations, and how does it compare to breeding populations? How does \Ag change during adaptation to a new environment, and can we detect signatures of additive channels in genomic data? How does the framework extend to non-Gaussian selection and non-normal trait distributions? These questions point toward both theoretical development and empirical investigation.

\subsection{Concluding remarks}

Additivity is not a property of genes but a property of populations in context. A population can behave additively on a curved landscape if its variance is confined to relatively flat regions. The additivity index \Ag provides a diagnostic for this condition, and the framework developed here connects that diagnostic to the dynamics of selection and inheritance. This perspective suggests that the success of additive models may reflect not the absence of gene interactions but the tendency of selection to confine populations to regions where those interactions contribute little to fitness variance. Whether this tendency is strong enough to explain observed patterns is an empirical question. We hope the framework developed here provides useful tools for addressing it.

% ==========================================================
\section*{Data Availability}
% ==========================================================
All code used to generate figures and run simulations is available at \url{https://github.com/dortizbarrientos/additive-channels} under an MIT licence. Parameter files and scripts required to reproduce all figures are included in the repository. No new empirical genotype or phenotype data were generated for this study.

% ==========================================================
\section*{Acknowledgments}
% ==========================================================
We thank colleagues in the ARC Centre of Excellence for Plant Success in Nature and Agriculture for helpful discussions during the development of this work. We thank, in particular, Bruce Walsh for insightful discussions about the origin and maintenance of quantitative genetic variation in nature. 

% ==========================================================
\section*{Conflicts of Interest}
% ==========================================================
The authors declare no competing interests.

\section*{Funding}
ARC Centre of Excellence for Plant Success in Nature and Agriculture [CE200100015].

\bibliographystyle{genetics}
\bibliography{references}

% ==========================================================
% APPENDICES
% ==========================================================
\appendix
% Bulletproof appendix-numbering machinery.
%
% Background: in some journal classes (gsajnl appears to be one),
% \appendix interacts with \section in ways that prevent the
% standard \renewcommand{\thesection}{S\arabic{section}} pattern
% from producing visible numbers. The symptom is that
% "Appendix~\ref{app:foo}" in the body renders as "Appendix " with
% a blank afterwards, even though the label exists in the .aux file.
%
% Fix: drive the numbering off our own counter (appsec) rather than
% the section counter, and use \section* with manual labelling so
% the heading does not depend on \appendix's behaviour at all.
% \refstepcounter sets \@currentlabel to \theappsec, which the
% subsequent \label captures, so \ref{app:foo} returns "S1", "S2",
% ..., and "Appendix~\ref{app:foo}" renders as "Appendix S1", "Appendix
% S2", etc. \phantomsection gives hyperref a proper anchor for the
% clickable link.
\newcounter{appsec}
\renewcommand{\theappsec}{S\arabic{appsec}}
\makeatletter
\newcommand{\appsection}[2]{%
  \refstepcounter{appsec}%
  \section*{Appendix \theappsec: #2}%
  \phantomsection%
  \label{#1}%
  \addcontentsline{toc}{section}{Appendix \theappsec: #2}%
}
\makeatother

\appsection{app:variance}{Variance Identities for Linear and Quadratic Forms}

This appendix derives the variance formulas used in the main text. We proceed step by step, aiming to be accessible to readers with a basic background in linear algebra. The main results are standard in multivariate analysis and matrix calculus; for a comprehensive reference on quadratic forms in normal variables see \citet{MathaiProvost1992}, and for the underlying matrix calculus see \citet{Anderson2003}, \citet{Mardia1979}, and \citet{MagnusNeudecker1988}.

\subsection{Mathematical tools}

The algebra uses standard results from multivariate analysis and matrix calculus. We cite canonical sources so readers can verify each step.

For a random vector $\vect{a}$ with covariance $\mat{G}$ and any fixed vector $\vect{b}$, $\Var(\vect{b}^\top\vect{a})=\vect{b}^\top\mat{G}\vect{b}$; see \citet{Anderson2003} or \citet{Mardia1979}. If $\vect{X}\sim\mathcal{N}(\vect{0},\Sigma)$ and $\mat{A}$ is symmetric, then $\Var(\vect{X}^\top\mat{A}\vect{X})=2\,\tr[(\mat{A}\Sigma)^2]$; see \citet{MagnusNeudecker1988} or \citet{Mardia1979}. Odd moments vanish for centred Gaussians, which is used to show $\Cov(L,Q)=0$ when $L$ is linear, and $Q$ is quadratic in $\vect{a}$; see \citet{Mardia1979}. We also use the cyclicity of traces and the diagonalisation of symmetric matrices; see \citet{HornJohnson2012}.

\subsection{Setup}

Let $\vect{a}$ be a random vector in $\RR^n$ with mean zero and covariance $\mat{G} = \E[\vect{a}\vect{a}^\top]$. For the quadratic variance result, we assume $\vect{a} \sim \mathcal{N}(\vect{0}, \mat{G})$ \citep{Anderson2003}.

\subsection{Variance of a linear form}

Let $L = \vect{b}^\top \vect{a}$ where $\vect{b}$ is a fixed vector. Since $\E[\vect{a}] = \vect{0}$, we have $\E[L] = 0$ and:
\begin{align}
\Var(L) &= \E[L^2] = \E[(\vect{b}^\top \vect{a})(\vect{a}^\top \vect{b})] \\
&= \E[\vect{b}^\top \vect{a}\vect{a}^\top \vect{b}] \\
&= \vect{b}^\top \E[\vect{a}\vect{a}^\top] \vect{b} \\
&= \vect{b}^\top \mat{G} \vect{b},
\end{align}
a standard identity for linear combinations of random vectors \citep{Anderson2003,Mardia1979}.

For $n=2$ with $\mat{G} = \text{diag}(\sigma_1^2, \sigma_2^2)$ and $\vect{b} = (b_1, b_2)^\top$, this gives $\vect{b}^\top \mat{G} \vect{b} = b_1^2 \sigma_1^2 + b_2^2 \sigma_2^2$.

\subsection{Variance of a quadratic form}

Let $Q = \tfrac{1}{2}\vect{a}^\top \mat{H} \vect{a}$ with $\vect{a} \sim \mathcal{N}(\vect{0}, \mat{G})$ and $\mat{H}$ symmetric.

A standard theorem for quadratic forms of centred Gaussians states that if $\vect{X} \sim \mathcal{N}(\vect{0},\Sigma)$ and $\mat{A}$ is symmetric, then
\begin{equation}
\Var(\vect{X}^\top \mat{A} \vect{X}) = 2\,\tr[(\mat{A}\Sigma)^2],
\end{equation}
see \citet{MagnusNeudecker1988} or \citet{Mardia1979}. Applying this with $\mat{A} = \tfrac{1}{2}\mat{H}$ and $\Sigma = \mat{G}$ gives
\begin{equation}
\Var(Q) = 2\,\tr\left[\left(\tfrac{1}{2}\mat{H}\mat{G}\right)^2\right] = \tfrac{1}{2}\,\tr[(\mat{H}\mat{G})^2],
\end{equation}
which is Equation~\eqref{eq:Vquad} in the main text.

\subsection{Covariance between linear and quadratic forms}

For $L = \vect{b}^\top\vect{a}$ and $Q = \tfrac{1}{2}\vect{a}^\top\mat{H}\vect{a}$,
\[
\Cov(L, Q) = \E[LQ] - \E[L]\E[Q] = \E[LQ],
\]
since $\E[L]=0$.

The quantity $\E[LQ]$ is a sum of third-order moments of $\vect{a}$. For a centred multivariate normal, all odd central moments vanish by symmetry \citep{Mardia1979}. Hence $\E[LQ] = 0$ and $\Cov(L, Q) = 0$ under the Gaussian assumption.

\subsection{Coordinate invariance of $\Vlin$ and $\Vquad$}

Let $\mat{D}$ be an invertible linear change of coordinates in breeding-value space, $\vect{a}\mapsto \mat{D}\vect{a}$. Standard chain-rule and covariance arguments give the induced transformations
\[
\mat{G}\mapsto \mat{D}\mat{G}\mat{D}^{\top}, \qquad \vect{b}\mapsto \mat{D}^{-\top}\vect{b}, \qquad \mat{H}\mapsto \mat{D}^{-\top}\mat{H}\mat{D}^{-1}.
\]
Direct substitution into the linear-variance functional yields
\[
\vect{b}^{\top}\mat{G}\vect{b}\;\mapsto\;(\mat{D}^{-\top}\vect{b})^{\top}(\mat{D}\mat{G}\mat{D}^{\top})(\mat{D}^{-\top}\vect{b})
=\vect{b}^{\top}\mat{G}\vect{b}.
\]
For the quadratic-variance functional, the matrix product $\mat{H}\mat{G}$ transforms by similarity,
\[
\mat{H}\mat{G}\;\mapsto\;\mat{D}^{-\top}\mat{H}\mat{G}\mat{D}^{\top},
\]
and using cyclic permutation of the trace,
\[
\tr[(\mat{H}\mat{G})^{2}]\;\mapsto\;\tr\!\left[(\mat{D}^{-\top}\mat{H}\mat{G}\mat{D}^{\top})^{2}\right]=\tr[(\mat{H}\mat{G})^{2}].
\]
Both $\Vlin$ and $\Vquad$ are therefore invariant under invertible linear changes of coordinates, and so is $\Ag$. This is the algebraic content of the basis-free expression in Equation~\eqref{eq:Ag_rotated}.

\appsection{app:selection}{Selection Operators}

This appendix derives the Gaussian selection operators used in the main text. Both derivations follow standard multivariate normal updating and matrix identities; see \citet{Anderson2003} for Gaussian conditioning and \citet{GolubVanLoan2013} for Woodbury-type identities.

\subsection{Level 1 derivation}

Prior: $\vect{a} \sim \mathcal{N}(\vect{0}, \mat{G})$ with density $p(\vect{a}) \propto \exp(-\tfrac{1}{2}\vect{a}^\top \mat{G}^{-1} \vect{a})$.

Fitness: $W(\vect{a}) \propto \exp(-\tfrac{1}{2}\vect{a}^\top \Gamma \vect{a})$.

Post-selection density:
\begin{align}
p(\vect{a} \mid \text{sel}) &\propto p(\vect{a}) \cdot W(\vect{a}) \\
&\propto \exp\left(-\tfrac{1}{2}\vect{a}^\top (\mat{G}^{-1} + \Gamma) \vect{a}\right).
\end{align}

This is Gaussian with precision $\mat{G}^{-1} + \Gamma$, hence covariance $\Gsel = (\mat{G}^{-1} + \Gamma)^{-1}$. This adding-precision interpretation is standard for Gaussian updates \citep{Anderson2003}.

\subsection{Level 2 derivation}

Let phenotype be $\vect{z} = \vect{a} + \vect{e}$ with $\vect{a} \sim \mathcal{N}(\vect{0}, \mat{G})$, $\vect{e} \sim \mathcal{N}(\vect{0}, \mat{E})$ independent, so $\mat{P} = \mat{G} + \mat{E}$. Gaussian stabilising selection with width $\Omega$ can be represented as a pseudo-observation model $\vect{y} = \vect{z} + \vect{\epsilon}$, $\vect{\epsilon} \sim \mathcal{N}(\vect{0}, \Omega)$, conditioned on $\vect{y} = \vect{0}$. Since $\vect{z} = \vect{a} + \vect{e}$, the effective observation noise is $\mat{E} + \Omega$.

Using the standard conditional covariance formula for a jointly Gaussian partition \citep{Anderson2003}, $\Cov(\vect{a} | \vect{y}) = \Sigma_{aa} - \Sigma_{ay}\Sigma_{yy}^{-1}\Sigma_{ya}$, and noting $\Sigma_{aa} = \mat{G}$, $\Sigma_{yy} = \mat{P} + \Omega$, and $\Sigma_{ay} = \mat{G}$, we obtain $\Gsel = \mat{G} - \mat{G}(\mat{P} + \Omega)^{-1}\mat{G}$.

\subsection{Convergence of Level 2 to Level 1}

When $\mat{E} = \mat{0}$, $\mat{P} = \mat{G}$, so $\Gsel = \mat{G} - \mat{G}(\mat{G} + \Omega)^{-1}\mat{G}$. Applying the Woodbury matrix identity (matrix inversion lemma) \citep{GolubVanLoan2013}, $\mat{G} - \mat{G}(\mat{G} + \Omega)^{-1}\mat{G} = (\mat{G}^{-1} + \Omega^{-1})^{-1}$. Setting $\Gamma = \Omega^{-1}$ recovers Level 1.

\subsection{The LD retention coefficient $\rho$: derivations under the Lande, Bulmer, and Barton--Turelli frameworks}
\label{app:rho_derivations}

The inheritance recursion in the main text (Eq.~\ref{eq:bulmer-LD}) treats LD retention as a single scalar~$\rho \in [0,1]$, where $\rho$ is the fraction of selection-induced gametic-phase disequilibrium retained in the next generation after recombination. This subsection provides theoretical foundations for the scalar~$\rho$ closure under three classical multilocus frameworks, identifies the conditions under which the closure is exact, and substantiates the claim that inbreeding raises~$\rho$ as an analytical rather than verbal argument.

\subsubsection{The Lande infinitesimal model}

In the infinitesimal limit, where additive variance arises from infinitely many freely recombining loci of infinitesimal effect, gametic-phase disequilibrium between any two loci decays geometrically. For two loci with recombination fraction~$r$, the standard random-mating recursion is $D_{ij}^{(t+1)} = (1-r)\, D_{ij}^{(t)}$ in the absence of selection \citep{Bulmer1971}. For unlinked loci, $r = 1/2$, so $D$ halves each generation.

The Bulmer reduction in additive variance under selection is built from a sum over locus-pair LD contributions \citep{Bulmer1971,Bulmer1980}. Writing $V_A^{(t,\mathrm{LD})}$ for the LD-induced component of additive variance in generation~$t$, the recursion for an unlinked-loci infinitesimal architecture under random mating is therefore
\begin{align}
V_A^{(t+1,\mathrm{LD})} &= (1-r)\, V_A^{(t,\mathrm{LD})} \;=\; \tfrac{1}{2}\, V_A^{(t,\mathrm{LD})}.
\label{eq:rho_lande}
\end{align}
Identifying the right-hand side with $\rho \cdot V_A^{(t,\mathrm{LD})}$ in the notation of Eq.~\ref{eq:bulmer-LD} gives $\rho = 1/2$ as the canonical value for the Lande infinitesimal model.

\subsubsection{The Bulmer model with arbitrary recombination architecture}

When loci are not all unlinked, the per-pair recombination fraction~$r_{ij}$ varies across pairs. The same per-pair LD recursion $D_{ij}^{(t+1)} = (1-r_{ij})\, D_{ij}^{(t)}$ applies, but now the variance recursion becomes a weighted sum. Following the structure of \citet{Bulmer1971} eq.~3.4 and writing $w_{ij}$ for the contribution of locus pair $(i,j)$ to the LD-induced component of additive variance,
\begin{align}
V_A^{(t+1,\mathrm{LD})} &= \frac{\sum_{i\neq j} w_{ij} (1 - r_{ij})\, D_{ij}^{(t)}}{\sum_{i\neq j} w_{ij}\, D_{ij}^{(t)}} \cdot V_A^{(t,\mathrm{LD})}.
\label{eq:rho_bulmer_general}
\end{align}
The leading factor is~$\rho$ for the Bulmer model under arbitrary recombination architecture. It collapses to a single scalar when the per-pair retention $1 - r_{ij}$ is approximately uniform across the pairs that contribute to the relevant element of~$\mat{G}$, in which case
\begin{align}
\rho_{\mathrm{Bulmer}} \;=\; \mathbb{E}_{ij}[\,1 - r_{ij}\,].
\label{eq:rho_bulmer_scalar}
\end{align}
For unlinked loci, $\rho_{\mathrm{Bulmer}} = 1/2$, recovering the Lande result. For tightly linked loci, $r_{ij} \to 0$ and $\rho_{\mathrm{Bulmer}} \to 1$. For a mixture of tight and loose linkages, $\rho$ takes intermediate values. This is the analytical content of the manuscript's verbal claim that linkage architecture moves~$\rho$ along the unit interval.

\subsubsection{Inbreeding and selfing raise $\rho$}

Under partial selfing at rate~$s$, recombination between loci has reduced effect on between-individual LD because selfed offspring are increasingly homozygous and recombination among identical haplotypes does not change pairwise LD at the population level. \citet{Wright1933} established the foundational result; \citet{HaymanMather1953} and \citet{Wright1969} give the multilocus extension. For unlinked loci ($r = 1/2$) and selfing rate $s \in [0, 1]$, the steady-state per-generation LD retention factor in a partially selfing population takes the form
\begin{align}
\rho(s) \;\approx\; \tfrac{1}{2}(1 + F),
\label{eq:rho_selfing}
\end{align}
where $F = s/(2 - s)$ is the inbreeding coefficient at selfing--outcrossing equilibrium \citep{Wright1969}. Two limits are immediate: $s = 0$ (random mating) gives $F = 0$ and $\rho = 1/2$, recovering the Lande value; $s = 1$ (complete selfing) gives $F = 1$ and $\rho = 1$, so all selection-induced LD persists across generations.

This is the formal version of the manuscript's argument that breeding programs achieve high~$\rho$ through inbreeding. Equation~\ref{eq:rho_selfing} translates the qualitative observation into a single numerical mapping between mating system and across-generation LD persistence: a breeding program at $s = 0.9$ has $F \approx 0.82$ and $\rho \approx 0.91$, deep into the regime where Bulmer-induced compression of~$\mat{G}$ persists across cycles, contrasting sharply with the $\rho = 0.5$ of randomly mating outcrossing populations.

\subsubsection{The Barton--Turelli multilocus framework}

\citet{BartonTurelli1991} derived the full multilocus recursion for the additive genetic variance under selection, recombination, and mutation, generalising Bulmer beyond pair-wise locus interactions to arbitrary numbers of loci with arbitrary recombination architecture. The key structural feature of their derivation is that the across-generation transmission of LD-induced variance reduction is governed by a recombination operator~$\mathcal{R}$ acting on the multilocus genotype distribution. The eigenvalues of~$\mathcal{R}$ index how disequilibria decay along distinct directions of the multilocus state space.

In this more general framework, the scalar~$\rho$ of Eq.~\ref{eq:bulmer-LD} emerges as a leading eigenvalue of~$\mathcal{R}$ under three conditions: (i)~weak selection, so that recombination rather than selection dominates the per-generation update of allele frequencies; (ii)~near-equilibrium allele frequencies, so that the recombination operator acts on a distribution that is close to its stationary form under random mating; and (iii)~symmetric recombination architecture, so that the direction-dependent decay rates collapse to a single scalar for the relevant element of~$\mat{G}$. These are the conditions under which the scalar closure is exact. When they fail, $\rho$ becomes a direction-dependent operator, an extension we identify as future work in the main text (the \textit{Inheritance and Equilibrium} section). \citet{Burger2000} provides a textbook treatment that situates the Barton--Turelli framework alongside the Lande and Bulmer formulations.

\subsubsection{Summary}

Across the three frameworks, the scalar~$\rho$ used in the manuscript's inheritance recursion has consistent interpretations: $\rho = 1/2$ in the Lande infinitesimal limit under random mating, $\rho = \mathbb{E}_{ij}[1 - r_{ij}]$ in the Bulmer model with non-uniform recombination, $\rho \approx (1+F)/2$ under partial selfing, and a leading eigenvalue of the recombination operator in the Barton--Turelli framework under symmetry conditions. The verbal argument that ``inbreeding raises~$\rho$'' translates into the analytical statement that $\rho$ is monotone increasing in the inbreeding coefficient~$F$, with $\rho \to 1$ as $F \to 1$. The verbal argument that ``tight linkage raises~$\rho$'' translates into the analytical statement that $\rho$ is monotone decreasing in the average inter-locus recombination fraction, with $\rho \to 1$ as $r \to 0$.

\appsection{app:equilibrium}{Equilibrium Analysis}

This appendix presents the algebra that leads to Equation~\eqref{eq:sigma-star} in the scalar isotropic case. The recursion is a simplified Bulmer-style inheritance update \citep{Bulmer1971, BartonTurelli1991}, combined with Gaussian stabilising selection (Level~1).

At equilibrium, the genetic variance satisfies
\[
\sigma^{*2} = \sigma^2_g + \rho\left(\frac{\sigma^{*2}}{1 + \gamma \sigma^{*2}} - \sigma^2_g\right),
\]
where $\sigma^2_g$ is the genic variance (which may itself reflect mutation input on the chosen timescale). Here $\gamma = \omega^{-2}$ is the selection intensity, where $\omega^2$ is the width of the Gaussian fitness function (see Level~1, the \textit{Level 1: Selection directly on breeding values} subsection).

Let $x = \sigma^{*2}$ and $s = \sigma^2_g$. Starting from
\[
x = s + \rho\left(\frac{x}{1+\gamma x} - s\right)
= s(1-\rho) + \frac{\rho x}{1+\gamma x},
\]
multiply both sides by $(1+\gamma x)$ to obtain
\[
x(1+\gamma x) = s(1-\rho)(1+\gamma x) + \rho x.
\]
Expanding and collecting terms yields the quadratic equation
\[
\gamma x^2 + (1-\rho)(1-\gamma s)\,x - s(1-\rho) = 0.
\]
Taking the positive root gives
\begin{equation}
\sigma^{*2}
= \frac{-(1-\rho)(1-\gamma \sigma^2_g)
+ \sqrt{(1-\rho)^2(1-\gamma \sigma^2_g)^2 + 4\gamma \sigma^2_g(1-\rho)}}{2\gamma}.
\label{eq:sigma-star-S3}
\end{equation}

This expression has the correct limiting behaviour: when $\rho = 0$ (no LD retention), $\sigma^{*2} = \sigma^2_g$; when $\rho \to 1$ (complete LD retention), $\sigma^{*2} \to 0$.

\appsection{app:example}{A Worked Example: Connecting Operations to Concepts}

This appendix uses a concrete numerical example to show how each mathematical operation connects to biological concepts. The goal is to let readers trace how abstract formulas become interpretable quantities. We use a two-dimensional example that is rich enough to show orientation effects but simple enough to compute by hand.

\subsection{The biological setup}

Consider a population of plants with two quantitative traits: flowering time ($z_1$) and plant height ($z_2$). These traits are influenced by breeding values $\vect{a} = (a_1, a_2)^\top$, where $a_1$ represents the additive genetic contribution to flowering time and $a_2$ represents the additive genetic contribution to height.

The population has the following genetic covariance matrix:
\begin{equation}
\mat{G} = \begin{pmatrix} 1.0 & 0.3 \\ 0.3 & 0.6 \end{pmatrix}
\end{equation}

Flowering time has additive genetic variance $G_{11} = 1.0$, height has variance $G_{22} = 0.6$, and the genetic correlation between them is $r_G = 0.3/\sqrt{1.0 \times 0.6} \approx 0.39$. This positive correlation might arise from pleiotropy (genes affecting both traits) or linkage disequilibrium between loci affecting different traits. The genetic correlation means selection on one trait induces a correlated response in the other; the direction depends on the sign of selection on the focal trait. In other words, because $ G_{12}$ is greater than zero, selection favours increased flowering time ($b_1>0$), which also tends to increase height as a correlated response.

\subsection{The fitness landscape}

Selection acts through a log-fitness surface with local geometry characterised by:
\begin{equation}
\vect{b} = \begin{pmatrix} 0.8 \\ 0.4 \end{pmatrix}, \quad \mat{H} = \begin{pmatrix} -0.4 & 0 \\ 0 & -0.4 \end{pmatrix}
\end{equation}

The gradient $\vect{b}$ tells us that fitness increases when $a_1$ increases (later flowering is favoured) and when $a_2$ increases (taller plants are favoured), with flowering time under stronger directional selection ($b_1 = 0.8 > b_2 = 0.4$). The Hessian $\mat{H}$ is negative definite, indicating stabilising selection around an optimum. Both traits experience equal curvature ($\gamma = 0.4$), and there is no correlational selection ($H_{12} = 0$). The curvature magnitude matrix is $\Gamma = -\mat{H} = 0.4\mat{I}$.

\subsection{Step 1: Computing $\Vlin$}

The linear variance formula is $\Vlin = \vect{b}^\top \mat{G} \vect{b}$. First, multiply $\mat{G}\vect{b}$:
\begin{equation}
\mat{G}\vect{b} = \begin{pmatrix} 1.0 & 0.3 \\ 0.3 & 0.6 \end{pmatrix} \begin{pmatrix} 0.8 \\ 0.4 \end{pmatrix} = \begin{pmatrix} 0.92 \\ 0.48 \end{pmatrix}
\end{equation}

The vector $\mat{G}\vect{b}$ is the expected response to selection---the direction and magnitude the population mean would move under the breeder's equation $\Delta\bar{\vect{a}} = \mat{G}\vect{\beta}$. The response is stronger at flowering time (0.92) than at height (0.48), reflecting both stronger selection on flowering time and greater genetic variance in that trait.

Completing the quadratic form:
\begin{equation}
\Vlin = \vect{b}^\top(\mat{G}\vect{b}) = (0.8, 0.4) \begin{pmatrix} 0.92 \\ 0.48 \end{pmatrix} = 0.928
\end{equation}

$\Vlin = 0.928$ is the additive genetic variance in local log-fitness under the linear approximation; it is closely related to the additive genetic variance in relative fitness that features in Fisher/Price formulations under weak selection.

\subsection{Step 2: Computing $\Vquad$}

The quadratic variance formula is $\Vquad = \tfrac{1}{2}\tr[(\mat{H}\mat{G})^2]$. This requires computing the matrix product $\mat{H}\mat{G}$, squaring it, and taking the trace.

\begin{equation}
\mat{H}\mat{G} = \begin{pmatrix} -0.4 & 0 \\ 0 & -0.4 \end{pmatrix} \begin{pmatrix} 1.0 & 0.3 \\ 0.3 & 0.6 \end{pmatrix} = \begin{pmatrix} -0.4 & -0.12 \\ -0.12 & -0.24 \end{pmatrix}
\end{equation}

The matrix $\mat{H}\mat{G}$ combines curvature with genetic variance. Each entry describes how curvature in one direction interacts with variance in another. The diagonal terms ($-0.4$ and $-0.24$) reflect the interaction between curvature and variance for the same trait; the off-diagonal terms ($-0.12$) reflect how genetic covariance couples the curvature effects.

\begin{equation}
(\mat{H}\mat{G})^2 = \begin{pmatrix} -0.4 & -0.12 \\ -0.12 & -0.24 \end{pmatrix}^2 = \begin{pmatrix} 0.1744 & 0.0768 \\ 0.0768 & 0.0720 \end{pmatrix}
\end{equation}

\begin{equation}
\tr[(\mat{H}\mat{G})^2] = 0.1744 + 0.0720 = 0.2464
\end{equation}
\begin{equation}
\Vquad = \tfrac{1}{2} \times 0.2464 = 0.1232
\end{equation}

$\Vquad = 0.1232$ captures curvature/nonlinearity in the local log-fitness surface with respect to breeding values; this can reflect epistasis, nonlinear phenotype–fitness mapping, or both. It is substantially smaller than $\Vlin$ because the population's genetic variance is modest relative to the curvature of the landscape.

\subsection{Step 3: Computing the additivity index}

\begin{equation}
\Ag = \frac{\Vlin}{\Vlin + \Vquad} = \frac{0.928}{0.928 + 0.1232} = \frac{0.928}{1.0512} \approx 0.883
\end{equation}

The additivity index $\Ag \approx 0.88$ means that 88\% of the genetic variance in log-fitness comes from the linear (additive) term, while only 12\% comes from curvature. This population is within an additive channel. Predictions based on the breeder's equation should work reasonably well; Fisher's theorem should be a good approximation; and estimated breeding values should accurately rank individuals by expected fitness.

\subsection{Step 4: What selection does to $\mat{G}$}
Using the Level 2 selection operator with environmental covariance $\mat{E} = 0.5\mat{I}$ and selection width $\Omega = 2.5\mat{I}$ (corresponding to $\Gamma = 0.4\mat{I}$):
\begin{equation}
\mat{P} = \mat{G} + \mat{E} = \begin{pmatrix} 1.5 & 0.3 \\ 0.3 & 1.1 \end{pmatrix}
\end{equation}
\begin{equation}
\mat{P} + \Omega = \begin{pmatrix} 4.0 & 0.3 \\ 0.3 & 3.6 \end{pmatrix}
\end{equation}
\begin{equation}
(\mat{P} + \Omega)^{-1} \approx \begin{pmatrix} 0.252 & -0.021 \\ -0.021 & 0.280 \end{pmatrix}
\end{equation}
\begin{equation}
\Gsel = \mat{G} - \mat{G}(\mat{P}+\Omega)^{-1}\mat{G} \approx \begin{pmatrix} 0.74 & 0.19 \\ 0.19 & 0.48 \end{pmatrix}
\end{equation}
Selection has compressed genetic variance. The variance in flowering time dropped from 1.0 to 0.74 (26\% reduction), and the variance in height dropped from 0.6 to 0.48 (19\% reduction). The genetic correlation changed slightly, from 0.39 to $0.19/\sqrt{0.74 \times 0.48} \approx 0.32$. This compression is the Bulmer effect in action: selection eliminates extreme genotypes, leaving a narrower distribution among survivors.

\subsection{Step 5: How compression affects \Ag}

After selection, we can recompute the additivity index using $\Gsel$:

\begin{equation}
\Vlin^{\text{sel}} = \vect{b}^\top \Gsel \vect{b} \approx 0.67
\end{equation}

For the quadratic variance:
\begin{equation}
\mat{H}\Gsel = \begin{pmatrix} -0.294 & -0.075 \\ -0.075 & -0.194 \end{pmatrix}
\end{equation}
\begin{equation}
\Vquad^{\text{sel}} = \tfrac{1}{2}\tr[(\mat{H}\Gsel)^2] \approx 0.068
\end{equation}

The new additivity index:
\begin{equation}
\Ag^{\text{sel}} = \frac{0.67}{0.67 + 0.068} \approx 0.91
\end{equation}

Selection increased \Ag from 0.88 to 0.91 in this single-generation calculation. This is the within-step self-correcting dynamic described in the main text: selection compresses variance, which reduces $\Vquad$ faster than $\Vlin$ (because $\Vquad \propto \sigma^4$ while $\Vlin \propto \sigma^2$), pushing the population deeper into the additive channel \emph{within the current generation}. Whether this gain is sustained across generations depends on whether $|\vect{b}|$ remains substantial. Under maintained directional selection (artificial selection or a moving optimum), $|\vect{b}|$ is held approximately constant by external pressure and \Ag continues to climb until variance reaches its compressed equilibrium. Under adaptation toward a fixed internal optimum, $|\vect{b}|$ collapses with the displacement on the same timescale as variance compresses, and \Ag falls toward zero. The single-step calculation here illustrates the mechanism; the \textit{Inheritance and Equilibrium} section and Figure~\ref{fig:adaptation-trajectory} show how the two regimes diverge over many generations.

\subsection{Summary}

This example illustrates the complete framework. Local geometry ($\vect{b}$, $\mat{H}$) describes the slope and curvature of the fitness surface that the population experiences. Genetic covariance ($\mat{G}$) describes which directions in breeding-value space the population actually explores. Variance decomposition ($\Vlin$, $\Vquad$) quantifies how much fitness variation comes from the linear versus quadratic terms of the Taylor expansion. The additivity index (\Ag) summarises whether the population is in a regime where additive prediction works. Selection operators ($\Gsel$) describe how selection compresses genetic variance within a generation. The self-correcting dynamic shows that variance compression raises \Ag, pushing populations toward additive channels under the conditions described in the main text.

The key insight is that additivity is not a property of the genes but of the population's variance relative to the local landscape geometry. A population with the same genetic architecture can be in an additive channel or outside it, depending on its variance and where it sits on the fitness surface. As discussed in the main text, this self-correcting dynamic can reverse when populations approach fitness optima where $|\vect{b}| \to 0$.

\appsection{app:simulation}{Stochastic-simulation validation of the additive-channel framework}

% Local figure-counter renumbering for the supplementary figures of this
% appendix.  We reset the figure counter and redefine \thefigure to
% display "S5.<n>" rather than the arabic main-text number, so the
% three figures here are numerically separated from the main-text
% figures (1--5).  Any later appendix that wants its own figure
% numbering should issue its own \setcounter{figure}{0} and
% \renewcommand{\thefigure}.  No subsequent appendix in this paper
% includes figures, so the counter need not be restored.
\setcounter{figure}{0}
\renewcommand{\thefigure}{S5.\arabic{figure}}

This appendix reports a simulation-based validation of the additive-channel framework using forward-time individual-based simulations in SLiM~\citep{HallerMesser2019}. The framework's central analytical claim is the identity \mbox{$\Ag = R^2$} between the analytical additivity index and the empirical fraction of log-fitness variance captured by linear regression on breeding values. We test this identity, and several derived predictions about the across-replicate reproducibility of the framework's geometric quantities, in finite populations evolving under sustained directional selection followed by stabilising selection at a new optimum. All code, parameters, and per-replicate output are available in the public repository (\href{https://github.com/dortizbarrientos/additive-channels}{github.com/dortizbarrientos/additive-channels}).

\subsection*{Simulation design}

We simulated an explicit-genome diploid population evolving on a multivariate fitness landscape. Population size was fixed at $N=2000$. The genome carried 200 quantitative-trait loci (QTL), 50 per trait, distributed across four traits. Per-locus effects on each trait were drawn from independent Gaussians with trait-specific standard deviations $\bm{\sigma} = (0.346, 0.283, 0.200, 0.141)$. The genotype-to-breeding-value map was strictly additive: each individual's breeding value on trait~$k$ was the sum of its allelic effects on trait~$k$ across all 50 loci of that trait. Selection was Gaussian-stabilising in phenotype space, with curvature (selection strength) $\bm{\gamma} = (0.10, 0.05, 0.04, 0.03)$ and a moving optimum on trait~1 only. The optimum drifted at velocity $0.3$ per generation for the first 40 generations, then froze at $\theta = (12, 0, 0, 0)$. Each replicate ran for 200 generations and reported the full breeding-value cloud at each generation. We used 20 replicates with sequential seeds $42,\ldots,61$.

Three features of this design merit comment. First, by displacing the optimum we ensure that $\vect{b}\neq\vect{0}$ during the directional phase, placing the population in the regime where \Ag is most informative; freezing the optimum at generation~40 then drives $\vect{b}\to\vect{0}$, allowing us to observe the framework's behaviour as it enters its known scope limit (the \textit{Additivity Index} section). Second, the linear genotype-to-breeding-value map separates our test from the literature on epistasis-mediated variance dynamics under nonlinear $G\!\to\!P$ maps~\citep{CarterHermissonHansen2005,Hansen2013}: any departure from monotone variance compression that we observe must arise from selection-induced linkage disequilibrium and drift acting on a strictly additive architecture, not from $G\!\to\!P$ nonlinearities. Third, the four-trait dimension lets us compute geometric invariants of $\mat{G}$ that would be trivial in two dimensions.

\subsection*{The identity \texorpdfstring{$\Ag = R^2$}{Ag = R\textasciicircum 2} holds across stochastic realisations}

At each generation of each replicate we computed three quantities from the breeding-value cloud and the Hessian $\mat{H}=-\bm{\gamma}$: the linear variance $\Vlin = \vect{b}^{\top}\mat{G}\vect{b}$, the quadratic variance $\Vquad = \tfrac{1}{2}\tr[(\mat{H}\mat{G})^2]$, and the analytical additivity index $\Ag = \Vlin/(\Vlin+\Vquad)$. We compared \Ag against the empirical $R^2$ from the Walsh decomposition of log-fitness on breeding values within each generation~\citep{Walsh2018}.

Figure~\ref{fig:s5_multiseed} shows the trajectory across 20 replicates. The analytical \Ag (mean across replicates, solid line) and empirical $R^2$ (mean, dashed line) are visually indistinguishable across the entire trajectory, including during the rapid rise (generations 1--40) when the population is tracking the moving optimum, the peak at the generation-40/41 transition, and the decay after the optimum freezes. Across all 200 generations and 20 replicates the median absolute difference $|\Ag - R^2|$ remains 0.013 to three decimal places, and the algebraic identity \mbox{$2\Vquad = \sum_i \mu_i^2$} holds to floating-point precision (maximum numerical discrepancy $<10^{-15}$). Two observations follow.

\paragraph{The identity is not specific to equilibrium.} 
The framework's analytical derivation is local and pointwise: it asserts equality between $\Ag$ and $R^2$ at any \mbox{$(\vect{b},\mat{H},\mat{G})$} configuration consistent with the local-Gaussian closure, regardless of whether the population is at, near, or far from any equilibrium. The simulation evaluates the identity at every generation of every replicate, providing an empirical test at hundreds of distinct off-equilibrium configurations along each trajectory. The identity holds throughout, including during phases when the local-Gaussian assumption is stressed by rapid mean displacement.

\paragraph{The identity passes through the framework's scope limit cleanly.} 
After generation~50, the population settles near the new optimum and the gradient $\vect{b}\to\vect{0}$. Both \Ag and $R^2$ collapse toward zero in lockstep, as the framework predicts: with $\Vlin \to 0$ but $\Vquad$ remaining strictly positive, the ratio $\Vlin/(\Vlin+\Vquad)\to 0$. This is not a failure of the identity but the framework's known scope limit (the \textit{Additivity Index} section). The simulation confirms that the limit is approached smoothly and identically by both sides of the analytical equation.

\begin{figure*}[ht]
\centering
\includegraphics[width=\textwidth]{multiseed_4trait_moving.pdf}
\caption{Multi-seed trajectory across 20 stochastic realisations. \textbf{Left:} the analytical $\Ag$ (solid line, mean of 20 replicates) and empirical $R^2$ (dashed line, mean) overlap throughout the trajectory, with the $5\text{--}95\%$ percentile band shown for $\Ag$. The grey trace shows replicate~1 individually. \textbf{Centre:} the four sorted eigenvalues of $\mat{G}$ across replicates, with $5\text{--}95\%$ percentile bands. \textbf{Right:} the gradient magnitude $|\vect{b}|$ rises during the directional phase (generations 1--40), peaks around the point at which the optimum stops moving, and then decays as the population catches up to the frozen optimum.}
\label{fig:s5_multiseed}
\end{figure*}

\subsection*{Across-replicate variability: the framework's invariants are tighter than their constituents}

A natural worry about the identity \mbox{$\Ag = R^2$} is that it might hold trivially in any individual replicate but be unstable across replicates---that is, that the absolute value of \Ag at a given generation might vary so much from seed to seed that the diagnostic is empirically uninformative. We tested this directly by computing the across-replicate coefficient of variation \mbox{$\mathrm{CV}(x) = \mathrm{sd}(x)/|\mathrm{mean}(x)|$} per generation for each metric.

Figure~\ref{fig:s5_cv} reports the result. At the peak (generation~40), \Ag has \mbox{$\mathrm{CV}=0.064$} and $R^2$ has \mbox{$\mathrm{CV}=0.056$}. The constituent quantities $\Vlin$, $\Vquad$, and $|\vect{b}|$ have CVs of $0.088$, $0.362$, and $0.174$, respectively. The individual eigenvalues of $\mat{G}$ have median CV $0.191$ and range $0.130$--$0.249$, whereas the trait-fixed diagonals $G_{kk}$ have median CV $\approx 0.237$. Two readings follow.

\paragraph{\Ag is approximately three times tighter across replicates than the individual eigenvalues of \texorpdfstring{$\mat{G}$}{G}.} 
This is the headline reproducibility result: at the peak, the framework's diagnostic is not merely identical to $R^2$ in each replicate, it is also more reproducible \emph{across} replicates than the underlying components. The mechanism is straightforward. \Ag is a scalar function of \mbox{$(\vect{b},\mat{H},\mat{G})$} that depends only on the trace-like quantities $\Vlin$ and $\Vquad$, not on the full eigenstructure of $\mat{G}$. Independent stochastic realisations of finite-population dynamics produce $\mat{G}$ matrices whose individual eigenvalues vary substantially (because which alleles fix in which order is genuinely seed-dependent), but whose trace-like summaries vary much less. \Ag inherits the tighter variability of the summaries rather than the looser variability of the eigenvalues.

\paragraph{The CV inflates after the population reaches the optimum, but for a known and benign reason.} 
After generation~50, the across-replicate CV of \Ag rises and eventually exceeds unity (Figure~\ref{fig:s5_cv}, left panel). This is not a failure of the framework but the same scope limit discussed above: with mean $\Ag \to 0$, the ratio $\mathrm{sd}/|\mathrm{mean}|$ inflates because of the vanishing denominator, not because of any qualitative change in the dynamics. The trait-fixed diagonals and the individual eigenvalues, which do not collapse to zero, retain stable CV in the range $0.20$--$0.25$ throughout. The diagnostic identity is most informative precisely where the framework was designed to be informative, namely the regime where $\vect{b}\neq\vect{0}$ and there is directional variance for the linear approximation to capture.

\begin{figure*}[ht]
\centering
\includegraphics[width=\textwidth]{additiveChannels/figures/cv_4trait_moving.pdf}
\caption{Across-replicate variability of framework quantities. \textbf{Left:} per-generation coefficient of variation, $\mathrm{CV}(x)=\mathrm{sd}(x)/|\bar{x}|$, across 20 replicates for $\Ag$, empirical $R^2$, $\Vlin$, $\Vquad$, $|\vect{b}|$, trait-fixed diagonals of $\mat{G}$, sorted eigenvalues of $\mat{G}$, and curvature-weighted eigenvalues $\mu_i$. The vertical dotted line marks generation~40, when the moving optimum freezes. The rise in $\mathrm{CV}(\Ag)$ and $\mathrm{CV}(R^2)$ after the population approaches the optimum reflects the known $\vect{b}\to\vect{0}$ scope limit, where the denominator in $\mathrm{sd}(x)/|\bar{x}|$ becomes small, not a loss of the identity $\Ag=R^2$. \textbf{Right:} CV at generation~40, showing $\Ag=0.064$, empirical $R^2=0.056$, median sorted-$\mat{G}$-eigenvalue CV $=0.191$, and median trait-diagonal CV $\approx 0.237$.}
\label{fig:s5_cv}
\end{figure*}

\subsection*{Geometric invariants of \texorpdfstring{$\mat{G}$}{G}: shape is reproducible even when axes are not}

The reproducibility result above raises a more pointed question: \emph{why} is \Ag tighter than its eigenvalue components? One hypothesis is that the geometric \emph{shape} of $\mat{G}$ (its effective dimensionality and anisotropy) is reproducible across replicates even when the individual axes are not. To test this we computed two scalar functions of $\mat{G}$'s eigenvalue spectrum that depend only on the \emph{set} of eigenvalues and not on any labelling:
\begin{align}
\keff(t) &\;=\; \frac{(\tr\mat{G}(t))^2}{\tr[\mat{G}(t)^2]} \;=\; \frac{\bigl(\sum_i \lambda_i\bigr)^2}{\sum_i \lambda_i^2}, 
\label{eq:s5_keff} \\[2pt]
\kappa(\mat{G}(t)) &\;=\; \frac{\lambda_{\max}(\mat{G}(t))}{\lambda_{\min}(\mat{G}(t))}. 
\label{eq:s5_cond}
\end{align}
The first quantity is the participation ratio (the \textit{Dimensional Considerations} section); it ranges in $[1, n_{\text{traits}}]$, with $\keff = n_{\text{traits}}$ when $\mat{G}$ is isotropic and $\keff = 1$ when $\mat{G}$ has collapsed to a single direction. The second is the condition number of $\mat{G}$ as a covariance matrix and quantifies how stretched the variance ellipsoid is. Both are immune by construction to the eigenvalue-relabelling that occurs when sorted eigenvalues swap order during the trajectory.

Figure~\ref{fig:s5_geom} reports the results. \keff describes a non-monotone trajectory: it rises from $2.95$ at generation~1 to a peak of $3.27$ at generation~36, then falls to a stable $2.74$ by generation~100. The condition number traces a complementary path: it falls from $5.91$ to a minimum near generation~40 (the most isotropic configuration $\mat{G}$ reaches), then rises steadily through generation~200 as the smallest eigenvalue continues to fall. Three statements follow.

\paragraph{The early rise of \keff is variance redistribution, not variance expansion.} 
The total genetic variance $\tr\mat{G}$ falls monotonically from $6.29$ at generation~1 to $3.32$ at generation~200; total variance is being lost throughout. The transient rise of \keff during generations 1--30 reflects \emph{redistribution} of the available variance among directions rather than any expansion of $\mat{G}$. Because our genotype-to-breeding-value map is strictly additive by construction, this dynamic is mechanistically distinct from the $\mat{G}$-expansion under directional selection that has been described in the presence of nonlinear $G\!\to\!P$ maps and positive epistasis~\citep{CarterHermissonHansen2005}; the latter cannot operate in our simulation. What we observe instead is a generic property of multivariate Gaussian-stabilising selection acting on an anisotropic $\mat{G}$: directions with the largest standing variance experience the largest absolute compression per generation, so the spectrum becomes transiently more uniform before the most-selected direction continues compressing past parity.

\paragraph{\keff is much more reproducible across replicates than $\kappa(\mat{G})$.} 
At the peak generation, $\mathrm{CV}(\keff) = 0.059$, comparable to $\mathrm{CV}(\Ag) = 0.064$ and well below the median individual-eigenvalue CV of $0.191$. By contrast $\mathrm{CV}(\kappa) = 0.240$, slightly higher than the eigenvalue median. The contrast is mathematically expected: \keff is built from the full sum and sum-of-squares of eigenvalues and is therefore robust to noise in any single eigenvalue, whereas $\kappa$ is a ratio of the largest to the smallest eigenvalue and is dominated by noise in $\lambda_{\min}$, which is itself the noisiest quantity in the spectrum. The geometric invariance result therefore has structure: \emph{not all sort-invariant scalars are equally reproducible}, and the framework's reliance on trace-like summaries rather than condition-number-like ratios is part of why \Ag inherits good reproducibility.

\paragraph{$\Ag$ and \keff are jointly reproducible but causally distinct.} 
At the peak, \Ag and \keff have similar CV ($0.064$ and $0.059$). It is tempting to read this as a causal claim that dimensional collapse is what drives the population into the additive channel, but the timing argues against that interpretation. \keff peaks five generations \emph{before} \Ag, and the rise of \Ag is dominated by the rise of $|\vect{b}|^2$ in the numerator $\Vlin = \vect{b}^{\top}\mat{G}\vect{b}$, not by changes in the geometric shape of $\mat{G}$. The right reading is that \Ag and \keff are \emph{both} sort-invariant geometric scalars of the local landscape, and \emph{both} inherit reproducibility from the same underlying mechanism (the variance-dampening of trace-like summaries under stochastic eigenvalue noise), but neither causes the other along the trajectory.

\begin{figure*}[ht]
\centering
\includegraphics[width=\textwidth]{geom_invariants_4trait_moving.pdf}
\caption{Geometric invariants of $\mat{G}$ across replicates. \textbf{Top-left:} effective dimension $\keff(t)$ across 20 replicates, with mean and $5\text{--}95\%$ percentile band; the dashed horizontal line marks the isotropic value $n_{\text{traits}}=4$. \textbf{Top-right:} anisotropy $\kappa(\mat{G})=\lambda_{\max}/\lambda_{\min}$, with mean and $5\text{--}95\%$ percentile band on a logarithmic axis. \textbf{Bottom:} CV at generation~40 for $\keff$, $\kappa(\mat{G})$, and the sorted eigenvalues of $\mat{G}$. $\keff$ has $\mathrm{CV}=0.059$, well below the median eigenvalue CV of $0.191$, whereas $\kappa(\mat{G})$ has $\mathrm{CV}=0.240$, reflecting its sensitivity to the smallest eigenvalue.}
\label{fig:s5_geom}
\end{figure*}

\subsection*{Synthesis and limitations}

Across 20 stochastic replicates of a four-trait moving-optimum scenario with a strictly additive genotype-to-breeding-value map, three results emerge. \emph{First}, the analytical identity $\Ag = R^2$ holds at every generation of every replicate, including off-equilibrium configurations during the directional phase and the approach to the new optimum. \emph{Second}, the framework's diagnostic \Ag is approximately three times tighter across replicates than the individual eigenvalues of $\mat{G}$, and about three to four times tighter than several of the underlying constituent quantities, especially the trait-fixed diagonals and $\Vquad$; the algebraic mechanism for this tightness is that \Ag depends on trace-like summaries rather than the full eigenstructure. \emph{Third}, the geometric invariants $\keff$ and $\kappa(\mat{G})$ have qualitatively different across-replicate behaviour: \keff is reproducible to roughly the same tolerance as \Ag, whereas $\kappa(\mat{G})$ is substantially more variable, in line with their differing algebraic dependence on the spectrum.

These results carry a suggestive but not yet rigorous implication for parallel evolution. If two populations independently encounter similar selection regimes, the framework predicts that they will converge on similar \emph{geometric} configurations of $(\vect{b},\mat{H},\mat{G})$ even when the underlying allelic architectures differ---paralleling the long-standing observation that populations adapting to similar environments often arrive at similar phenotypic outcomes via different molecular routes~\citep{Schluter1996,Stern2013}. Our simulation tests a strong version of this claim, in which initial conditions are identical across replicates and only the random seed varies; whether the invariance generalises to populations differing in initial allele frequencies or mutational input is a question for future work, including the comparative-evolution treatment we develop in subsequent papers in this series.

Two notes on the scope of this validation. First, the simulation operates entirely at the breeding-value level. The genotype-to-breeding-value map is strictly additive by construction, and the variance dynamics we observe in $\mat{G}$ arise from two BV-level mechanisms---selection-induced linkage disequilibrium under Gaussian stabilising selection (the Bulmer effect, with subsequent erosion by recombination) and drift acting on allele frequencies. Both compress $\mat{G}$, consistent with the framework's prediction. A separate literature \citep{CheverudRoutman1995,Hansen2006,Hansen2013,CarterHermissonHansen2005} describes a different mechanism by which additive genetic variance can change under directional selection: when the underlying $G\!\to\!P$ map is itself nonlinear, allele-frequency shifts can alter the average effects of substitutions, converting variance between additive, dominance, and epistatic components, and under positive epistasis can transiently \emph{expand} $V_A$ rather than compress it. We do not test this mechanism here, by design rather than by omission. The framework treats breeding values as primitives, takes the BV distribution as given, and analyses the local geometry of log-fitness over BV-space (the \textit{Local Geometry of Log-Fitness Surfaces} section); whatever the upstream $G\!\to\!P$ map---linear or nonlinear, additive or richly epistatic---it delivers a distribution of breeding values on which the framework's diagnostic $\Ag$ then operates. The simulation reported here tests the BV-level part of this picture cleanly. Whether the framework's predictions remain accurate when $V_A$ is changing through allele-frequency-induced reorganisation rather than through BV-level compression alone is an upstream question, requiring simulations that couple a richer $G\!\to\!P$ map to the same diagnostic; we treat it as a follow-up problem.

Second, the simulation is at a fixed population size with a fixed genetic architecture. A full sensitivity analysis across $N$, the number of QTL per trait, the strength of stabilising selection, and the rate of optimum displacement would map the regime structure of the framework more completely; here we have aimed to validate the central identity and characterise reproducibility under a single non-trivial scenario rather than to chart the full regime space.
\appsection{app:dictionary}{Notation Dictionary Derivations}

This appendix derives the two identities used in the Lande--Arnold dictionary (the \textit{Notation: relation to Lande--Arnold (1983)} subsection): the Phillips--Arnold relation between log-fitness curvature and the relative-fitness quadratic gradient, and the environmental-buffering relation between curvature in breeding-value space and curvature in phenotype space.

\subsection{The relation $\mat{H} = \bm{\gamma}-\bm{\beta}\bm{\beta}^{\top}$}

Let $W$ be a twice-differentiable fitness function in phenotype space, and let $w(\vect{z}) = W(\vect{z})/\bar{W}$ be relative fitness. At the population mean phenotype $\bar{\vect{z}}$, relative fitness satisfies $w(\bar{\vect{z}})=1$. The Lande--Arnold quadratic gradient $\bm{\gamma}$ is the matrix of second partial derivatives of $w$ at $\bar{\vect{z}}$,
\[
\gamma_{ij} \;=\; \frac{\partial^2 w}{\partial z_i\,\partial z_j}\bigg|_{\bar{\vect{z}}},
\]
and the linear gradient is $\beta_i = (\partial w/\partial z_i)|_{\bar{\vect{z}}}$. Direct differentiation of $\log w$ gives
\[
\frac{\partial^2 \log w}{\partial z_i \partial z_j}
\;=\;
\frac{1}{w}\frac{\partial^2 w}{\partial z_i \partial z_j}
\;-\;
\frac{1}{w^2}\frac{\partial w}{\partial z_i}\frac{\partial w}{\partial z_j}.
\]
Evaluating at $\bar{\vect{z}}$ where $w=1$,
\[
\frac{\partial^2 \log w}{\partial z_i \partial z_j}\bigg|_{\bar{\vect{z}}}
\;=\;
\gamma_{ij} - \beta_i\beta_j.
\]
Since $\log w = \log W - \log\bar{W}$ differs from $\log W$ only by a constant, the Hessians of $\log w$ and $\log W$ coincide, and we recover Equation~\eqref{eq:H_vs_gamma}: $\mat{H} = \bm{\gamma}-\bm{\beta}\bm{\beta}^{\top}$ at the population mean. This is the relation given by \citet{PhillipsArnold1989}.

\subsection{Environmental buffering: $\mat{H}_{\vect{a}}$ versus $\mat{H}_{\vect{z}}$}

Consider Gaussian stabilising selection in phenotype space with optimum $\bm{\theta}$ and selection-width matrix $\bm{\Omega}$, so that $W(\vect{z}) \propto \exp\!\big(\!-\tfrac{1}{2}(\vect{z}-\bm{\theta})^{\top}\bm{\Omega}^{-1}(\vect{z}-\bm{\theta})\big)$. The phenotype-space curvature is
\[
\mat{H}_{\vect{z}} \;=\; \nabla^2_{\vect{z}}\log W(\vect{z}) \;=\; -\bm{\Omega}^{-1}.
\]
With phenotype $\vect{z} = \vect{a}+\vect{e}$ and environmental deviation $\vect{e}\sim\mathcal{N}(\vect{0},\mat{E})$ independent of $\vect{a}$, the breeding-value-space fitness function is
\[
W(\vect{a}) \;=\; \mathbb{E}_{\vect{e}}[W(\vect{a}+\vect{e})]
\;=\; \int \exp\!\big(\!-\tfrac{1}{2}(\vect{a}+\vect{e}-\bm{\theta})^{\top}\bm{\Omega}^{-1}(\vect{a}+\vect{e}-\bm{\theta})\big)\,\phi_{\mat{E}}(\vect{e})\,\mathrm{d}\vect{e},
\]
where $\phi_{\mat{E}}$ is the density of $\mathcal{N}(\vect{0},\mat{E})$. Completing the square in $\vect{e}$ and integrating gives a Gaussian in $\vect{a}$,
\[
W(\vect{a}) \;\propto\; \exp\!\Big(\!-\tfrac{1}{2}(\vect{a}-\bm{\theta})^{\top}(\bm{\Omega}+\mat{E})^{-1}(\vect{a}-\bm{\theta})\Big),
\]
so
\[
\mat{H}_{\vect{a}} \;=\; \nabla^2_{\vect{a}}\log W(\vect{a}) \;=\; -(\bm{\Omega}+\mat{E})^{-1}.
\]
Because $\bm{\Omega}+\mat{E}\succeq\bm{\Omega}$ (with strict inequality whenever $\mat{E}\succ\mat{0}$), the eigenvalues of $|\mat{H}_{\vect{a}}|$ are uniformly smaller than those of $|\mat{H}_{\vect{z}}|$: environmental variance buffers the apparent curvature. This is the same buffering that appears in the Level~2 selection update on $\mat{G}$ (the \textit{Selection Reshapes Genetic Covariances} section, Equation~\eqref{eq:level2}); the two appearances are different consequences of the same marginalisation in Equation~\eqref{eq:Wa_definition}.

\subsection{Putting the dictionary to work}

To compute \Ag from a published Lande--Arnold study, the workflow is: (i) start with the reported $(\bm{\beta},\bm{\gamma})$ together with $\mat{G}$, $\mat{E}$, and $\mat{P}$; (ii) form the phenotype-space log-fitness Hessian $\mat{H}_{\vect{z}} = \bm{\gamma}-\bm{\beta}\bm{\beta}^{\top}$; (iii) recover the BV-space Hessian via the buffering relation, $\mat{H}_{\vect{a}}^{-1} = \mat{H}_{\vect{z}}^{-1} - \mat{E}$, when $\mat{E}$ is small enough that this remains negative-definite; (iv) substitute into $\Vlin = \vect{b}^{\top}\mat{G}\vect{b}$ and $\Vquad = \tfrac{1}{2}\tr[(\mat{H}_{\vect{a}}\mat{G})^2]$, using $\vect{b}\approx\bm{\beta}$ to leading order \emph{under weak selection and after mapping both gradients to the same domain} (footnote~\ref{fn:b-vs-beta}). Outside that regime, $\bm{\beta}$ and $\vect{b}$ are formally distinct and the substitution should be replaced by the explicit BV-space gradient, computed from $\ell(\vect{a})$ directly when an analytic fitness function is available. The resulting \Ag is the diagnostic for the population at hand. When the buffering correction in (iii) cannot be made (e.g., when $\mat{E}$ is unknown, or when the system is far enough from the Gaussian regime that the closed-form mapping fails), one can still compute the phenotype-space analogue of \Ag using $\mat{H}_{\vect{z}}$ in place of $\mat{H}_{\vect{a}}$, recognising that this measures predictability of relative fitness from phenotype rather than predictability of log fitness from breeding value.

\appsection{app:lineage}{Historical Context}

This note provides a selective historical context for the framework in the main text. It is not a comprehensive review. The goal is to locate the additivity index \Ag and the ``additive channel'' interpretation within a well-established line of work on (i)~Gaussian approximations for polygenic traits, (ii)~the maintenance and structure of $\mat{G}$ under stabilising selection, and (iii)~the long-standing empirical observation that additive models often predict well despite widespread molecular interactions.

\paragraph{Selection-induced linkage disequilibrium.}
A first pillar is the realisation that selection can reshape observable additive variance without requiring immediate changes in allele frequencies. \citet{Bulmer1971} formalised how selection generates linkage disequilibrium (LD) that typically reduces additive genetic variance within a generation, with recombination eroding that LD across generations. This mechanism is one of the main ways in which $\mat{G}$ can be compressed on short timescales. It motivates our use of an effective LD-retention parameter $\rho$ as a coarse summary of how strongly selection-induced LD persists.

\paragraph{Maintenance of the G-matrix.}
A second pillar concerns why a $\mat{G}$-matrix exists and why it has structure in natural populations. \citet{Lande1980} analysed the maintenance of genetic variances and covariances under pleiotropic mutation and stabilising selection. Under Gaussian approximations, mutation introduces a multivariate input of variation, while stabilising curvature removes variation more strongly in directions that incur larger fitness costs. This provides a mechanistic route by which $\mat{G}$ can become structured under mutation--selection balance and helps justify treating $\mat{G}$ as an object with interpretable geometry rather than merely an empirical summary.

\paragraph{Local landscape measurement.}
A third pillar is the development of methods for measuring selection locally. \citet{LandeArnold1983} provided an operational framework for estimating linear and quadratic selection components from data, separating slope from curvature in a local approximation. Our notation differs---we work in breeding-value space and use $\vect{b}$ and $\mat{H}$ for log-fitness slope and curvature---but the conceptual move is the same: treat selection as locally described by first and second derivatives, then ask what those derivatives imply for response.

\paragraph{Limits of Gaussian closure.}
A fourth strand is the recognition that Gaussian closure can fail and that departures from normality matter. The ``house-of-cards'' approach \citep{Turelli1984} highlighted regimes in which mutation--selection balance does not yield a simple Gaussian picture, motivating caution about when orthogonality results (such as $\Cov(L,Q)=0$) can be expected to hold.

A useful refinement is that ``non-Gaussian'' often means \emph{leptokurtic} rather than merely skewed: in house-of-cards regimes, segregating variance can be dominated by rare alleles of relatively large effect, so real genetic variance can become more ``event-driven'' in finite populations even when long-run expectations remain bounded \citep{Turelli1984,KeightleyHill1988,JohnsonBarton2005}. This has two implications for our framework. First, orthogonality results that are exact under Gaussianity (e.g.\ $\Cov(L,Q)=0$) need not hold, so \Ag should be interpreted as a diagnostic tied to the Gaussian/quadratic closure rather 
than a universal identity. Second, stabilising selection and partial linkage can still generate negative LD and compress observable additive variance, but over longer time scales, changes in variance are often dominated by the allele-frequency dynamics of the few loci that contribute most variance, rather than by departures from normality alone \citep{TurelliBarton1994}. For this 
reason, we treat Gaussianity as a working approximation, and emphasise simulations and diagnostics that reveal when fat-tailed variation undermines the clean variance decomposition while leaving the qualitative ``variance-compression'' intuition intact.

\paragraph{Orientation, not just magnitude.}
A fifth line of development is the geometric emphasis that the orientation of genetic variance matters, not just its magnitude. \citet{Schluter1996} popularised the idea that long-term evolution is channelled along the major axis of genetic variation ($\vect{g}_{\max}$), and subsequent work has clarified how the geometry of variance and selection shapes evolutionary response. Our framework is compatible with this view but asks a different question: given a current neighbourhood of the fitness surface, does the population's variance cloud sample that neighbourhood in a way that makes linear prediction adequate? The diagnostic \Ag is designed to quantify this directly, and to track how it changes as selection compresses variance or as mutation and gene flow expand it.

\paragraph{Modern complex-trait syntheses.}
Finally, modern syntheses in complex-trait genetics emphasise that additive statistical models can succeed even when molecular interactions are common. \citet{HillGoddardVisscher2008} reviewed why additive variance often dominates in practice, and \citet{SellaBarton2019} discussed complex-trait evolution in a high-dimensional setting where many directions are effectively weakly selected.  The additive-channel framework complements these perspectives by providing a geometric diagnostic for when local linearisation is expected to be safe, and by making explicit the roles of variance compression (selection, LD persistence) and variance expansion (mutation, gene flow) in moving populations into or out of high-\Ag regimes.

\appsection{app:faq}{Anticipated Questions}

This appendix addresses questions that readers may raise about the scope, assumptions, and novelty of the framework. We frame these as frequently asked questions (FAQs) to make the responses easier to access.

\paragraph{Q1: Is this just a reformulation of known results?}

The individual components of this framework---the Bulmer effect \citep{Bulmer1971}, Lande--Arnold selection gradients \citep{LandeArnold1983}, variance decomposition under Gaussian assumptions \citep{Mardia1979}---are indeed well established. The contribution lies in their \emph{synthesis} into a real-time diagnostic that connects local landscape geometry to prediction accuracy. Prior work established that selection compresses variance and that fitness surfaces have local slope and curvature; the additivity index \Ag provides a quantitative bridge between these observations. Specifically, the result that $\Ag = \Corr(L, \tilde{\ell})^2$ (Equation~\ref{eq:Ag_as_R2}) gives the diagnostic direct operational meaning: \Ag is the $R^2$ of the best linear predictor for local log-fitness. This connection between a variance ratio and prediction accuracy was not, to our knowledge, made explicit in prior treatments.

\paragraph{Q2: Are the Gaussian assumptions too restrictive?}

The Gaussian assumption is a working approximation, not a claim about biological reality. We use it because it yields tractable variance identities and clean orthogonality ($\Cov(L,Q) = 0$). When breeding-value distributions deviate substantially from normality---as they do under house-of-cards mutation or strong selection on few loci \citep{Turelli1984, TurelliBarton1994}---the exact identities break down. However, our simulations (Figure~\ref{fig:Ag_validation}) show that even under heavy-tailed distributions, \Ag remains a monotone predictor of rank correlation between linear predictions and true fitness. The diagnostic degrades gracefully: it stops being an exact variance fraction but continues to track ``how much curvature scrambles linear ranking.'' We view \Ag as most informative when distributions are approximately Gaussian, and as a rough guide otherwise.

\paragraph{Q3: How would anyone measure \Ag empirically?}

Estimating \Ag requires three quantities: the genetic covariance matrix $\mat{G}$, the fitness gradient $\vect{b}$, and the fitness Hessian $\mat{H}$. The first two are routinely estimated in quantitative genetics; $\mat{G}$ from pedigree or genomic data, and $\vect{b}$ from selection differentials or regression of fitness on breeding values \citep{LandeArnold1983}. The Hessian is more challenging but not impossible: quadratic regression of fitness on phenotypes or breeding values provides estimates of curvature \citep{PhillipsArnold1989, BlowsBrooks2003}. Deep mutational scanning and combinatorial mutagenesis now provide direct measurements of local fitness landscapes, from which $\vect{b}$ and $\mat{H}$ can in principle be extracted. Even rough estimates would allow testing the framework's core prediction: that populations with higher \Ag show better alignment between additive predictions and observed responses.

\paragraph{Q4: What about allele frequency change?}

The recursion equation we presented in the main text is a phenomenological closure for variance dynamics, not a multilocus population-genetic model. We do not track allele frequencies, linkage disequilibrium at individual loci, or the effects of selection on specific variants. This is a deliberate scope limitation. The question we address is: given the current $(\mat{G}, \vect{b}, \mat{H})$, how does the geometry of variance relative to curvature affect prediction accuracy, and how does selection reshape that geometry? Answering this question does not require tracking which alleles are responsible for the variance. A full treatment would connect the variance-level dynamics studied here to allele-frequency dynamics, but that connection is beyond the present scope.

\paragraph{Q5: Does \Ag apply near fitness optima?}

At a fitness optimum where $\vect{b} = \vect{0}$, the linear variance $\Vlin = 0$ and hence $\Ag = 0$ by convention. This does not mean the population is ``dominated by epistasis''---it simply means there is no directional selection to compare against curvature. The additive-channel concept is most informative when populations are displaced from optima: tracking a moving target, experiencing directional selection, or adapting to a new environment. Near equilibrium under stabilising selection, \Ag may be low because the gradient is weak, not because curvature is strong. This property is a feature of additivity, not a bug: \Ag diagnoses when linear \emph{directional} prediction works, which is precisely the regime where the breeder's equation is applied.

\paragraph{Q6: Why work in log-fitness rather than fitness?}

Three reasons motivate the choice of log-fitness $\ell = \log W$. First, log-fitness is the natural scale for Gaussian selection models: if phenotypes are normal and fitness is Gaussian, the log-fitness surface is quadratic. Second, multiplicative fitness effects become additive on the log scale, aligning with biological intuitions about independent contributions to survival or reproduction \citep{Rice2004}. Third, working on the log scale avoids complications when fitness values approach zero. The cost is that our results apply to log-fitness variance, not fitness variance directly; the two are related but not identical.

\paragraph{Q7: How does dimensionality affect the framework?}

The curvature relevance ratio $R = \Vquad/\Vlin$ scales with effective dimensionality. In the isotropic case, $R \propto n\gamma^2\sigma^2/|\vect{b}|^2$, so higher dimensionality increases curvature's relative contribution for fixed variance and gradient. However, most genetic variance in high-dimensional systems is concentrated in a low-dimensional subspace \citep{Blows2007, Kirkpatrick2009}, so effective dimensionality $\keff$ may be much smaller than nominal dimensionality $n$. Whether natural populations occupy high-\Ag regimes in high-dimensional trait space remains an open empirical question. The framework provides tools for addressing it, but does not presume the answer.

\paragraph{Q8: This framework is in breeding-value space; how does it map to genotype landscapes?}

We use a local continuous coordinate system $\vect{z}$ because it is the natural tangent-space approximation to a discrete genotype neighbourhood. In quantitative genetics, $\vect{z}$ can be interpreted as breeding values (or any heritable trait coordinates) with population covariance $\mat{G}=\mathrm{Cov}(\vect{z})$. In a genotype landscape, one can instead take $\vect{z}$ to be any local genotype coordinate system (e.g., indicator variables for a set of mutations, allele-frequency deviations around a reference state, or other local encodings), and then approximate log-fitness by a second-order expansion around that neighbourhood. In that setting, $\vect{b}$ and $\mat{H}$ correspond to the local first- and second-derivative structure of log-fitness in those genotype coordinates, and epistasis manifests as curvature through $\mat{H}$. This mapping hides two assumptions: first, that a local neighbourhood exists where a quadratic approximation is adequate; second, that the coordinate system used to define $\mat{G}$ matches the genetic variation actually sampled by the population (and by the measurement design). Deep mutational scanning and combinatorial mutagenesis provide precisely the kind of local-neighbourhood data. In contrast, in natural populations $\vect{b}$ and $\mat{H}$ are typically inferred in trait or breeding-value space.

\paragraph{Q9: Is $V_{\mathrm{quad}}=\tfrac{1}{2}\mathrm{tr}\!\left[(\mat{H}\mat{G})^{2}\right]$ well-behaved if $\mat{H}\mat{G}$ is not symmetric?}

Yes. $V_{\mathrm{quad}}$ is defined as the variance contribution of the quadratic term $Q=\tfrac{1}{2}\vect{z}^{\mathsf{T}}\mat{H}\vect{z}$ under $\vect{z}\sim\mathcal{N}(\vect{0},\mat{G})$, with $\mat{G}$ symmetric positive semidefinite and $\mat{H}$ symmetric (as a Hessian). The trace form can be written without any symmetry requirement on $\mat{H}\mat{G}$:
\[
\mathrm{tr}(\mat{H}\mat{G}\mat{H}\mat{G})
=
\mathrm{tr}\!\left[\big(\mat{G}^{1/2}\mat{H}\mat{G}^{1/2}\big)^2\right].
\]
Because $\mat{G}^{1/2}\mat{H}\mat{G}^{1/2}$ is symmetric, the right-hand side equals $\sum_i \lambda_i^2\ge 0$, making non-negativity explicit and avoiding any need to assume that $\mat{H}$ and $\mat{G}$ commute.

\paragraph{Q10: Does $\Ag$ predict response error, or only log-fitness predictability?}

By construction, $\Ag$ is a predictability diagnostic for local log-fitness. Under the Gaussian/quadratic closure, $\Ag=\mathrm{Corr}(L,\tilde{\ell})^2$ is the $R^2$ of the best linear predictor of $\tilde{\ell}$ from $L$, so the residual log-fitness variance is exactly $V_{\mathrm{quad}}$. Equivalently, the root-mean-squared log-fitness error of the linear approximation scales as $\sqrt{V_{\mathrm{quad}}}$. Writing the curvature relevance ratio as $R=V_{\mathrm{quad}}/V_{\mathrm{lin}}$, we have $\Ag=1/(1+R)\approx 1-R$ when curvature is a small perturbation. Translating that perturbation into an explicit error for a particular response statistic (mean shift, ranking error, realised response magnitude) depends on the selection operator and on the distribution of $\vect{z}$. In that sense, $\Ag$ diagnoses when linear directional prediction is expected to be reliable, rather than fully determining response error on its own.

\paragraph{Q11: The inheritance recursion is phenomenological; why should we believe it?}

We agree that the recursion is a closure, not a complete multilocus population-genetic model (cf.\ Q4). Its purpose is to capture, with one parameter, a well-established mechanism: selection generates linkage disequilibrium that reduces additive variance (the Bulmer effect), and recombination erodes that disequilibrium over subsequent generations \citep{Bulmer1971}. In the classic random-mating, effectively unlinked setting, the LD generated in one generation decays on the order of one half in the next, corresponding to a retention factor around $\rho\approx 1/2$, whereas tight linkage, reduced recombination, selfing, or clonal reproduction imply higher retention (larger $\rho$). We therefore treat $\rho$ as an empirical knob summarising LD persistence for the trait basis used in $\mat{G}$. It can be estimated from observed post-selection changes in additive variance, from marker-based LD decay, or from repeated $\mat{G}$ estimates across generations. The recursion is used to study how variance compression and replenishment move populations into or out of high-\Ag regimes, not to replace detailed allele-frequency modeling.

\paragraph{Q12: Aren't thresholds like ``$\mathcal{A}_g>0.8$'' arbitrary?}

Yes. Any hard cutoff is a convenience rather than a biological law. The additivity index $\mathcal{A}_g$ is best treated as a \emph{continuous} diagnostic of local linear predictability, not as evidence for discrete ``phases'' of genetic architecture. The term ``additive channel'' is therefore descriptive: it denotes regions of parameter space where curvature contributes little to local log-fitness variance and where linear ranking is expected to work well. In practice, the appropriate threshold depends on the decision problem. If the goal is accurate \emph{ranking} of genotypes, a moderate $\mathcal{A}_g$ may suffice, since rank correlations can remain high even when $\mathcal{A}_g$ is not close to 1. If the goal is accurate \emph{prediction of response magnitude} or variance components, a higher $\mathcal{A}_g$ may be required. We therefore use illustrative values (e.g.\ $0.8$ or $0.9$) only to aid interpretation and visualization, and recommend reporting $\mathcal{A}_g$ as a continuous quantity (ideally with uncertainty) rather than relying on any universal cutoff.

\paragraph{Q13: How do the dynamics of the additive channel differ between natural populations and elite breeding programs?}
The behaviour of \Ag is governed by the balance between forces that contract genetic variance (stabilising selection and selection-induced LD, especially under high $\rho$) and forces that expand it (mutation, recombination, migration). In breeding programs, negligible mutational input on the relevant timescale and high LD retention through inbreeding or selfing drive populations deep into the additive channel; selection compresses $\mat{G}$ within generations and inheritance fails to restore much of that variance, so $\Vquad$ collapses faster than $\Vlin$ and \Ag tends toward $1$. In outcrossing natural populations, recombination erodes selection-induced LD each generation and mutation supplies fresh variance, so $\mat{G}$ stabilises at a mutation--selection balance and \Ag often settles at intermediate values. The two regimes are therefore not different theories but different parameter settings of the same framework, with the contrast captured by Figure~\ref{fig:trajectories} in the main text.

\end{document}